\documentclass[11pt]{amsart}
\usepackage{amssymb}
\usepackage{booktabs}
\usepackage[dvipsnames]{xcolor}
\usepackage{microtype}
\usepackage[colorlinks]{hyperref}
\hypersetup{linkcolor=MidnightBlue, citecolor=MidnightBlue, urlcolor=MidnightBlue}

\linespread{1.06}

% Long typewriter identifiers make justified lines hard to break; raise the
% breaking tolerance and give TeX extra stretch before it resorts to
% overfull lines.
\tolerance=3000
\emergencystretch=3em

\usepackage{listings}
\lstdefinelanguage{lean}{
  keywords={theorem,def,structure,abbrev,lemma,noncomputable,where,fun,let,by,Prop,Type,extends,instance,class},
  sensitive=true,
  comment=[l]{--},
  morecomment=[s]{/-}{-/},
  string=[b]",
}
\lstdefinestyle{lean}{
  language=lean,
  basicstyle=\ttfamily\footnotesize,
  keywordstyle=\bfseries\color{MidnightBlue},
  commentstyle=\itshape\color{gray},
  columns=fullflexible,
  keepspaces=true,
  breaklines=true,
  breakatwhitespace=true,
  frame=single,
  rulecolor=\color{gray!45},
  backgroundcolor=\color{gray!6},
  xleftmargin=2mm,
  aboveskip=3.5mm,
  belowskip=3.5mm,
  extendedchars=true,
  literate={¬}{{$\neg$}}1 {¹}{{${}^{1}$}}1 {ö}{{\"o}}1 {Λ}{{$\Lambda$}}1
    {Ω}{{$\Omega$}}1 {α}{{$\alpha$}}1 {β}{{$\beta$}}1 {λ}{{$\lambda$}}1
    {μ}{{$\mu$}}1 {σ}{{$\sigma$}}1 {φ}{{$\varphi$}}1 {ε}{{$\varepsilon$}}1
    {ᵐ}{{${}^{m}$}}1 {‖}{{$\|$}}1 {•}{{$\bullet$}}1
    {⁻¹}{{${}^{-1}$}}2 {⁻}{{${}^{-}$}}1 {₀}{{${}_{0}$}}1 {₁}{{${}_{1}$}}1
    {₂}{{${}_{2}$}}1 {ℂ}{{$\mathbb{C}$}}1 {ℕ}{{$\mathbb{N}$}}1
    {ℝ}{{$\mathbb{R}$}}1 {→}{{$\to$}}1 {↔}{{$\leftrightarrow$}}1
    {∀}{{$\forall$}}1 {∂}{{$\partial$}}1 {∃}{{$\exists$}}1 {∈}{{$\in$}}1
    {∉}{{$\notin$}}1 {∑}{{$\sum$}}1 {∞}{{$\infty$}}1 {∧}{{$\wedge$}}1
    {∫}{{$\int$}}1 {≠}{{$\neq$}}1 {≤}{{$\leq$}}1 {≥}{{$\geq$}}1
    {⊆}{{$\subseteq$}}1 {⟨}{{$\langle$}}1 {⟩}{{$\rangle$}}1
    {𝓝}{{$\mathcal{N}$}}1 {Θ}{{$\Theta$}}1 {∩}{{$\cap$}}1 {⇑}{{$\uparrow$}}1,
}


\subjclass[2020]{43A99, 43A85, 30H25, 42B15, 42B35, 42C15, 42C40, 28C99}
\keywords{atomic decomposition, Besov space, harmonic analysis, wavelets, multipliers, Haar wavelets, atoms, formal verification, Lean}

\title[Besov-ish spaces: the Lean 4 formalization]{Besov-ish spaces through atomic decomposition:\\ the Lean 4 formalization}

\author[D. Smania]{Daniel Smania}
\address{Departamento de Matem\'atica \\
  Instituto de Ci\^encias Matem\'aticas e de Computa\c{c}\~ao-Universidade de S\~ao Paulo (ICMC/USP) - S\~ao Carlos \\
  Caixa Postal 668 \\ S\~ao Carlos-SP \\ CEP 13560-970 \\ Brazil.}
\email{smania@icmc.usp.br}
\urladdr{\url{https://sites.icmc.usp.br/smania/}}

\newtheorem{theorem}{Theorem}[section]
\newtheorem{corollary}[theorem]{Corollary}
\newtheorem{lemma}[theorem]{Lemma}
\newtheorem{proposition}[theorem]{Proposition}
\theoremstyle{definition}
\newtheorem{definition}[theorem]{Definition}
\newtheorem{remark}[theorem]{Remark}

\newcommand{\SA}{{\mathcal A}}
\newcommand{\SB}{{\mathcal B}}
\newcommand{\SF}{{\mathcal F}}
\newcommand{\SG}{{\mathcal G}}
\newcommand{\SH}{{\mathcal H}}
\newcommand{\SP}{{\mathcal P}}
\newcommand{\SW}{{\mathcal W}}
% Lean identifiers and file paths are often very long; rendering them with
% a plain \texttt would overflow the line. \lean therefore inserts an
% invisible breakpoint before the name and after every separator `_`, `.`
% and `/`, so long names can wrap at natural places.
\ExplSyntaxOn
\tl_new:N \l__lean_tl
\NewDocumentCommand{\leanname}{m}
 {
  \tl_set:Nn \l__lean_tl {#1}
  \tl_replace_all:Nnn \l__lean_tl {\_} {\_\allowbreak}
  \tl_replace_all:Nnn \l__lean_tl {.} {.\allowbreak}
  \tl_replace_all:Nnn \l__lean_tl {/} {/\allowbreak}
  {\ttfamily \tl_use:N \l__lean_tl}
 }
\ExplSyntaxOff
% In math mode the name is set in an unbreakable \mbox (rare and short);
% in text mode an invisible leading breakpoint lets the whole name move to
% the next line when needed.
\newcommand{\lean}[1]{\ifmmode\mbox{\leanname{#1}}\else\hspace{0pt}\leanname{#1}\fi}

% In-text formalization notes (replacing the former footnotes): each
% statement is followed by a short paragraph naming the Lean declarations
% that formalize it and describing what each of them proves.
\newcommand{\formalpar}[1]{\par\addvspace{2.5mm}{\small\sloppy\noindent
  \textbf{\color{MidnightBlue}Formalization.}\ #1\par}\addvspace{2.5mm}}

% Box opening each section: the Lean file documented by the section and the
% project files it imports (its dependencies inside the repository).
% The stretchable glue on both sides of the box provides a good breakpoint
% after amsart's run-in subsection headings, so the box drops to its own
% line instead of overflowing next to the heading.
\newcommand{\filebox}[2]{\par\addvspace{3mm}
  {\noindent\hspace{0pt plus 1fill}\penalty0
  \fcolorbox{gray!55}{gray!6}{\parbox{0.955\linewidth}{\small\raggedright
  \textbf{File:} \lean{#1}\\
  \textbf{Imports (within the project):} \lean{#2}}}%
  \hspace{0pt plus 1fill}\par}\addvspace{3mm}}

\hypersetup{colorlinks=true, pdftitle={Besov-ish spaces through atomic decomposition: the results formalized in Lean}, pdfauthor={Daniel Smania}}

\begin{document}

\begin{abstract}
This document describes, without proofs, exactly those results of the paper
\emph{Besov-ish spaces through atomic decomposition} (Analysis \& PDE 15
(2022), no.\ 1) that have been formally verified in the Lean~4 repository
\texttt{BesovSpacesGoodGrid}, together with a few results obtained in the
formalization that go beyond the paper. It was produced from a prerelease
version of the published paper
(\lean{preprint-source/besovish-2020-07-18.tex} in the repository), but all
references to numbered statements of the paper
follow the \emph{published} version. Unlike the paper, the document is
organized to mirror the repository: Part~I covers the weak-grid library
(\texttt{WeakGrid/}) and Part~II the good-grid library (\texttt{GoodGrid/}),
with sections ordered by file dependency; each section opens with a box
naming the Lean file it documents and the project files that file imports.
Statements follow the notation of the paper; each result is followed, in the
running text, by a \emph{Formalization} paragraph naming the corresponding
Lean declarations, describing what each of them proves, and recording every
deviation from the published statements, together with a snapshot of the
Lean code. The document also explains the purpose and architecture of the
formalization and ends with a guide to the repository files. The aggregate
root of the repository compiles with no \texttt{sorry}, and the main
theorems depend only on the standard axioms \texttt{propext},
\texttt{Classical.choice}, \texttt{Quot.sound}; the only declared
\texttt{sorry} of the repository is one inner sublemma (the $u_1+u_2$
construction) of the in-progress file formalizing Pointwise
Multipliers~II, which is not yet imported by the root.
\end{abstract}

\maketitle

\setcounter{tocdepth}{2}
\tableofcontents

\section{Introduction}

\subsection{Purpose of the formalization}

The repository \texttt{BesovSpacesGoodGrid} formalizes, in Lean~4 over
mathlib, a substantial part of the atomic decomposition theory of
\cite{paper}. The goals are twofold. First, to obtain machine-checked proofs
of the main results of the paper --- the Banach structure, compactness and
completeness of Besov-ish spaces, the transmutation-of-atoms machinery, the
equivalence of the atomic, Haar, and mean-oscillation norms, and the theory of
pointwise multipliers up to the non-Archimedean property and its
positive-cone refinements. Second, and just as important, to produce a
\emph{reusable} Lean library: the basic theory is stated and proved for
abstract \emph{weak grids} and abstract atom families, and the good-grid /
Souza-atom spaces of the paper are then obtained as a specialization. The
weak-grid layer (grids, atoms, Besov-ish spaces, scales, completeness,
transmutation) is independent of the good-grid geometry and is intended to be
usable in later developments.

This document is organized to mirror the repository rather than the paper.
Part~I documents the weak-grid library
(\texttt{BesovSpacesGoodGrid/WeakGrid/}) and Part~II the good-grid library
(\texttt{BesovSpacesGoodGrid/GoodGrid/}); inside each part the sections
follow the import order of the Lean files, so that every section only uses
material from earlier sections. Each section opens with a box naming the
file it documents and the project files it imports.
The mathematical pipeline, in this order, is:
\begin{enumerate}
  \item define quantitative good grids and general weak grids;
  \item package weak-grid cells and local spaces, and define atom families
    with support, convexity, phase invariance, and atom-size estimates;
  \item build levelwise atomic blocks and infinite $L^p$ representations,
    define the Besov-ish space and the coefficient-cost gauge
    \lean{Norm\_Costpq};
  \item prove the $L^t$ embedding, scale inclusions, representation-limit
    theorems, compactness statements, and completeness for the cost norm;
  \item prove the transmutation claims moving atomic representations from one
    weak grid to another, and specialize them to compare Souza's atoms, Besov's
    atoms, and sandwiched atom families on good grids;
  \item introduce the positive cone, the Haar and standard representation
    gauges, and the mean-oscillation norm, with finite-constant comparison
    theorems;
  \item prove the Dirac-approximation identities: the partial sums of the
    standard representation evaluate as averages of the function against the
    grid Dirac kernels, with the resulting $L^\infty$ bounds on ancestor
    coefficient sums (paper Proposition~17.1);
  \item develop the pointwise-multiplier layer: the \lean{selfs} classes,
    boundedness of \lean{selfs} multipliers, the $\mathcal{B}^s_{1,1}$
    multiplier characterization, the non-Archimedean multiplier estimates,
    including their positive-cone versions, and their consequences --- the
    inclusion of the tail \lean{selfs} classes in the multiplier space,
    strongly regular domains and Pointwise Multipliers~I, and the
    in-progress Pointwise Multipliers~II.
\end{enumerate}
A guide to the repository files is given in Section~\ref{guide}.

\subsection{How this document was produced}

This is a mathematical description of the current state (June 11, 2026) of the
Lean~4 formalization of \cite{paper}. It was produced from a \emph{prerelease
version of the published paper} (the LaTeX source
\lean{preprint-source/besovish-2020-07-18.tex}, included in the repository)
by the following rules.

\begin{itemize}
\item Results that are not present in the repository were deleted. These are:
  the elementary estimates of paper Section~4 (H\"older-like and convolution
  tricks, inlined in the Lean proofs without standalone statements); the
  H\"older and bounded variation atoms (paper Sections~11B--11C) and the
  corresponding characterizations (paper Propositions~16.2 and~16.3); and all
  of paper Sections~19--21 (quasialgebra, description of $\SB^s_{1,1}$ by
  indicators, left compositions). Dirac's approximations (paper Section~17),
  Corollary~18.6, strongly regular domains and Pointwise Multipliers~I
  (paper 18.7--18.9) are formalized and documented in
  Sections~\ref{diracsec} and~\ref{multsec}; Pointwise Multipliers~II
  (paper 18.10) is stated, and its proof is complete except for one inner
  sublemma, the only declared \lean{sorry} of the repository (see
  Section~\ref{multsec}).
\item Statements were updated to match the precision of the formalization;
  every change with mathematical content is recorded in the
  \emph{Formalization} paragraph that follows the statement.
\item All references of the form ``paper Proposition~$x.y$'' use the
  numbering of the \emph{published} version, Analysis \& PDE \textbf{15}
  (2022), no.~1 (which the reader should consult alongside this document);
  the prerelease source in the repository has the same section structure but
  its statement numbers may differ.
\item The material was reordered to follow the repository (see above), so
  the section numbers of this document do \emph{not} match the section
  numbers of the paper.
\item No proofs are included. Each statement is followed by a
  \emph{Formalization} paragraph giving the corresponding Lean declarations
  and files, with a short description of what each declaration proves, and
  then by a snapshot of the Lean code that formalizes it (the formal
  statement as it appears in the repository, with the proof omitted; an
  ellipsis \texttt{...} marks the omitted proof or, in a few long
  definitions, omitted structural fields).
\end{itemize}

A global change must be singled out: \emph{throughout this document
$p \in [1,\infty)$, $q \in [1,\infty]$ and $u \in [1,\infty]$}. The paper
develops Part~I for the quasi-Banach range $p,q \in (0,\infty]$; the
formalization restricts to the Banach range, so $\rho=\min\{1,p,q\}=1$
everywhere and all the $\rho$-norms of the paper are genuine norms. In Lean
the standing hypotheses are \lean{1 <= p}, \lean{p != \(\infty\)},
\lean{Fact (1 <= q)}, \lean{1 <= u}, \lean{0 < s}.

\section{Notation}

\begin{table}[h]
  \centering
  \caption{Most important notation/symbols/constants}
  \label{tab:table1}
  \begin{tabular}{cc}
    \toprule
    Symbol & Description\\
    \midrule
    $I$ & phase space\\
    $m$ & finite measure in the phase space $I$\\
    $a_P, b_P$ & an atom supported on $P$\\
    $\mathcal{A}$ & a class of atoms \\
    $\mathcal{A}(Q)$ & set of atoms of class $\mathcal{A}$ supported on $Q$ \\
    $\mathcal{A}^{sz}_{s,p}$& class of $(s,p)$-Souza's atoms \\
    $\mathcal{A}^{bs}_{s,\beta,p,q}$ & class of $(s,\beta,p,q)$-Besov's atoms \\
    $\mathcal{P}$ & grid of $I$\\
    $\mathcal{P}^k$ & family of subsets of $I$ at the $k$-th level of $\mathcal{P}$\\
    $\lambda_1 \leq \lambda_2$ & describes geometry of the good grid $\mathcal{P}$\\
    $\mathcal{B}^s_{p,q}(\mathcal{A})$ & $(s,p,q)$-Banach space defined by the class of atoms $\mathcal{A}$\\
    $\mathcal{B}^s_{p,q}$ or $\mathcal{B}^s_{p,q}(\mathcal{A}^{sz}_{s,p})$ & $(s,p,q)$-Banach space defined by Souza's atoms\\
    $P, Q, W$ & elements of the grid $\mathcal{P}$\\
    $L^p$ or $L^p(m)$ & Lebesgue spaces of $(I,\mathbb{A}, m)$ \\
    $|\cdot|_p$ & norm in $L^p$, $p \in (0,\infty]$ \\
    $p'$ & $1/p+1/p'=1$, with $p\in [1,\infty]$ \\
    \bottomrule
  \end{tabular}
\end{table}

We removed from the paper's table the rows for H\"older atoms, bounded
variation atoms, the constants $\rho$ and $\hat t$ (no longer needed in the
Banach range) since the corresponding material is not formalized. Constants
denoted $C$, $C(\cdot)$ below depend only on the displayed arguments; the Lean
statements carry explicit values for most of them.

\vspace{8mm}
\centerline{\Large\bf I. THE WEAK-GRID LIBRARY}
\smallskip
\centerline{\large\texttt{BesovSpacesGoodGrid/WeakGrid/}}
\vspace{6mm}

\begin{center}
\fbox{\begin{minipage}{0.88\textwidth}
\centering
In Part I we assume $s>0$, \ $p \in [1,\infty)$, \ $u\in[1,\infty]$ \ and \
$q \in [1,\infty]$.\footnote{Paper: $s>0$, $p\in (0,\infty)$,
$q\in(0,\infty]$. See the Introduction for the restriction to the Banach
range.}
\end{minipage}}
\end{center}
\vspace{4mm}

\section{Measure spaces and weak grids}\label{partition}

\filebox{WeakGrid/Definition.lean}{none (external:
\textnormal{\texttt{UnbalancedHaarWavelet}},
\textnormal{\texttt{LaminarFamiliesMaximalBinaryTrees}})}

The starting point of the whole library is deliberately poor in structure:
no metric, no topology, no nesting between scales --- only a finite measure
space together with a sequence of finite families of measurable sets
(``cells''), one family per scale $k$, satisfying a single quantitative
hypothesis: a uniform bound $C_1$ on how many cells of the \emph{same} level
can meet a given cell. This bounded-overlap condition $G_1$ is exactly what
is needed to make levelwise $\ell^p$ bookkeeping work: when a level-$k$ sum
of locally supported functions is integrated, each point is covered by at
most $C_1$ summands, so $L^t$ norms of levelwise sums can be compared with
$\ell^p$ norms of their coefficients at the cost of powers of $C_1$ (this is
the computation behind inequality (6-5) of the paper). Everything in Part~I
--- the Besov-ish spaces, their completeness, compactness, and the
transmutation machinery --- is built on this hypothesis alone.

Let $I$ be a measure space with a $\sigma$-algebra $\mathbb{A}$ and $m$ be a
measure on $(I,\mathbb{A})$, $m(I)<\infty$. Given a measurable set $J\subset I$
denote $|J|=m(J)$. We denote the Lebesgue spaces of $(I,\mathbb{A},m)$ by
$L^p$. A {\bf weak grid} is a sequence of finite families of measurable sets with
positive measure $\mathcal{P}=(\mathcal{P}^k)_{k\in \mathbb{N}}$, so that at
least one of these families is non empty and
\begin{itemize}
\item[$G_1$.] Given $Q\in \mathcal{P}^k$, let
$$\Omega_Q^k=\{ P \in \mathcal{P}^k\colon \ P\cap Q\neq \emptyset \}.$$
Then
$$C_{1}=\sup_k \sup_{Q \in \mathcal{P}^k} \# \Omega_Q^k< \infty.$$
\end{itemize}
Define $|\mathcal{P}^k|=\sup \{|Q|\colon Q \in \mathcal{P}^k\}$. We often abuse
notation using $\mathcal{P}$ for both $(\mathcal{P}^k)_{k\in \mathbb{N}}$ and
$\cup_k \mathcal{P}^k$.

\formalpar{A grid is the structure \lean{WeakGrid} in
\lean{WeakGrid/Definition.lean} (fields \lean{partitions},
\lean{measurable}, \lean{positive\_measure}, \lean{exists\_nonempty},
\lean{Cmult1}, \lean{overlap\_card\_le}), bundled with its ambient measurable
space as \lean{WeakGridSpace}. The paper calls this simply a ``grid''; in the
repository it is called a \emph{weak grid}, to distinguish it from the good
grids of Section~\ref{goodgrids}. The paper's convention that $P \neq Q$ for
cells at different levels is replaced by indexing cells by (level, cell)
pairs (the structure \lean{WeakGridCell}).   Implementation notes: each level is a
\lean{Finset (Set \(\alpha\))}, so all levelwise sums are finite sums and no
convergence issues arise inside a level; the exponents $p,q,u$ live in
$\mathbb{R}_{\geq 0} \cup \{\infty\}$ (the type \lean{ENNReal}), which lets the
endpoint $q=\infty$ or $u=\infty$ be handled by the same statements instead
of by parallel ``usual modification'' versions; and the overlap neighborhood
$\Omega_Q^k$ is the definition \lean{overlapFinset}, on which $G_1$ is
imposed as a cardinality bound.}
\begin{lstlisting}[style=lean]
structure WeakGrid where
        /-- The measure on `α`. -/
        μ : MeasureTheory.Measure α
        /-- The measure is finite on the whole space. -/
        isFinite : MeasureTheory.IsFiniteMeasure μ
        /-- The finite family `P^k` at each level `k`. -/
        partitions : ℕ → Finset (Set α)
        /-- Every cell is measurable. -/
        measurable : ∀ k Q, Q ∈ partitions k → MeasurableSet Q
        /-- Every cell has positive measure. -/
        positive_measure : ∀ k Q, Q ∈ partitions k → 0 < μ Q
        /-- At least one level is nonempty. -/
        exists_nonempty : ∃ k, (partitions k).Nonempty
        /-- Uniform same-level overlap multiplicity bound. -/
        Cmult1 : ℕ
        /-- `# Ω_Q^k ≤ Cmult1` for every `Q ∈ P^k`. -/
        overlap_card_le :
          ∀ k Q, Q ∈ partitions k → (overlapFinset (partitions k) Q).card ≤ Cmult1

structure WeakGridSpace where
        grid : WeakGrid (α := α)
\end{lstlisting}

\begin{remark}
There are plenty of examples of spaces with grids; the simplest one is
$[0,1)$ with the Lebesgue measure and the dyadic grid
$\mathcal{D}^k=\{[i/2^k,(i+1)/2^k),\ 0\leq i< 2^k\}$, and similarly the dyadic
and $d$-adic grids of $[0,1)^n$. Any measure space with a finite non-atomic
measure can be endowed with a grid made of a nested sequence of partitions
whose elements at the same level have precisely the same measure.
\end{remark}

\section{Atoms}\label{secatoms}

\filebox{WeakGrid/Atoms.lean}{WeakGrid/Definition.lean}

An atom family assigns to every cell $Q$ a set $\mathcal{A}(Q)$ of model
functions supported on $Q$, normalized so that an atom of
a cell of measure $|Q|$ has size at most  $|Q|^{s-1/(u'p)}$ in $L^{pu}$. The
normalization is chosen so that the smoothness parameter $s$ plays, on an
abstract measure space, the role that derivative order plays on
$\mathbb{R}^n$: smaller cells must carry atoms with larger  norm before
they contribute unit cost, exactly as in the classical atomic
decompositions of Besov spaces. Two structural axioms deserve comment.
Convexity and phase invariance of $\mathcal{A}(Q)$ ($A_3$) are what make
the infimum norm of Section~\ref{besovish} a genuine norm: given
representations of $f$ and of $g$, the triangle inequality is proved by
merging them cellwise, and the merged ``atom'' is a convex combination of
phase-rotated atoms, hence again an atom (this is the proof of paper
Proposition~6.3). The optional compactness axiom $A_5$ is what later turns
boundedness in cost into compactness in $L^p$. The flexibility of the
abstraction --- a possibly different local model space $\mathcal{B}(Q)$ for
every cell --- is used in earnest in Part~II, where the same interface is
instantiated by one-dimensional spaces of constants (Souza's atoms) and by
infinite-dimensional induced Besov spaces (Besov's atoms).

Let $\mathcal{P}$ be a grid. Let $p \in [1,\infty)$, $u \in [1,\infty]$, and
$s>0$. A family of {\bf atoms} associated with $\mathcal{P}$ of type $(s,p,u)$
is an indexed family $\mathcal{A}$ of pairs
$(\mathcal{B}(Q),\mathcal{A}(Q))_{Q\in \cup_k \mathcal{P}^k}$ where
\begin{itemize}
\item[$A_1$.] $\mathcal{B}(Q)$ is a complex linear space of functions
contained in $L^{pu}$.
\item[$A_2$.] If $\phi \in \mathcal{B}(Q)$ then $\phi(x)=0$ for every
$x \not\in Q$.
\item[$A_3$.] $\mathcal{A}(Q)$ is a nonempty convex subset of $\mathcal{B}(Q)$
such that $\phi \in \mathcal{A}(Q)$ implies
$\sigma\phi\in\mathcal{A}(Q)$ for every $\sigma \in \mathbb{C}$ satisfying
$|\sigma|=1$.
\item[$A_4$.] We have $$|\phi|_{pu} \leq |Q|^{s-\frac{1}{u'p}}$$ for every
$\phi \in \mathcal{A}(Q)$, where $1/u + 1/u' = 1$.
\end{itemize}
We will say that $\phi \in \mathcal{A}(Q)$ is an $\mathcal{A}$-atom of type
$(s,p,u)$ supported on $Q$. Sometimes we also assume
\begin{itemize}
\item[$A_5$.] For every $Q\in \mathcal{P}$ we have that $\mathcal{A}(Q)$ is a
compact subset in the strong topology of $L^p$,
\end{itemize}
or
\begin{itemize}
\item[$A_6$.] every $\mathcal{A}(Q)$ is a sequentially compact subset in the
weak topology of $L^p$.
\end{itemize}

\formalpar{Atom families are the structure \lean{AtomFamily} in
\lean{WeakGrid/Atoms.lean}, whose fields formalize $A_1$--$A_4$. The local
space $\mathcal{B}(Q)$ is the structure \lean{LocalVectorSpace} (field
\lean{localSpace}): in the paper $A_1$ asks for a complex \emph{Banach} space
contained in $L^{pu}$, while Lean only requires a complex vector space
together with an injective linear map into the functions
$\alpha \to \mathbb{C}$, all whose members belong to $L^{pu}$ (field
\lean{local\_memLp}); no completeness of the local space is assumed, and it
is never used. $A_2$ is the field \lean{local\_support}; $A_3$ corresponds
to \lean{atoms\_nonempty}, \lean{atoms\_convex} and
\lean{atoms\_phase\_invariant} (the paper states the phase invariance as an
``if and only if''; the two formulations are equivalent); $A_4$ is the field
\lean{atom\_bound}. The compactness assumption $A_5$ appears as the
predicate \lean{AssumptionA5}. The paper's alternative assumption $A_7$
($\mathcal{B}(Q)$ finite dimensional and $\mathcal{A}(Q)$ closed) is used in
the paper only to derive $A_5$; in the formalization the finite-dimensional
case is instantiated directly for Souza's atoms. Implementation notes: the
H\"older conjugate $u'$ is a structure field (\lean{uConj}) constrained by
the mathlib-style predicate \lean{ENNReal.HolderConjugate}, rather than a
computed value, which avoids case splits at $u \in \{1,\infty\}$; the local
space construction \lean{LocalVectorSpace} is universe-polymorphic, so the
same \lean{AtomFamily} interface accepts both small concrete models (the
constants $\mathbb{C}$, for Souza's atoms) and large function-space models
(induced $L^p$ spaces, for Besov's atoms); and the atom-size bound is
stated with \lean{MeasureTheory.eLpNorm}, the
$\mathbb{R}_{\geq0}\cup\{\infty\}$-valued $L^p$ seminorm of mathlib, so no
integrability side conditions are needed to state $A_4$.}
\begin{lstlisting}[style=lean,basicstyle=\ttfamily\scriptsize]
structure AtomFamily
  (G : WeakGridSpace (α := α)) (s : ℝ) (p u : ℝ≥0∞) where
  /-- The Hölder conjugate exponent `u'`. -/
  uConj : ℝ≥0∞
  /-- The smoothness parameter is positive. -/
  s_pos : 0 < s
  /-- `p ∈ [1,∞)`. -/
  one_le_p : 1 ≤ p
  /-- `p < ∞`, expressed in `ℝ≥0∞` as `p ≠ ∞`. -/
  p_ne_top : p ≠ ∞
  /-- `u ∈ [1,∞]`. -/
  one_le_u : 1 ≤ u
  /-- `uConj` is the Hölder conjugate of `u`. -/
  holder_conjugate : ENNReal.HolderConjugate u uConj
  /-- The local vector space `B(Q)`. -/
  localSpace : WeakGridCell G → LocalVectorSpace.{u, v} α
  /-- The chosen atoms `A(Q)`, as elements of `B(Q)`. -/
  atoms : ∀ Q, Set ((localSpace Q).carrier)
  /-- Every cell has at least one atom. -/
  atoms_nonempty : ∀ Q, (atoms Q).Nonempty
  /-- `B(Q)` is contained in `L^{pu}`. -/
  local_memLp : ∀ Q φ, MeasureTheory.MemLp ((localSpace Q).toFun φ) (p * u) G.measure
  /-- Local functions are supported on `Q`. -/
  local_support : ∀ Q φ, ∀ x, x ∉ Q.cell → (localSpace Q).toFun φ x = 0
  /-- `A(Q)` is convex in the local vector space. -/
  atoms_convex : ∀ Q, Convex ℝ (atoms Q)
  /-- `A(Q)` is invariant under multiplication by complex scalars of modulus one. -/
  atoms_phase_invariant :
    ∀ Q φ (σ : ℂ), φ ∈ atoms Q → ‖σ‖ = (1 : ℝ) → σ • φ ∈ atoms Q
  /-- The atom size estimate in `L^{pu}`. -/
  atom_bound :
    ∀ Q φ, φ ∈ atoms Q →
      MeasureTheory.eLpNorm ((localSpace Q).toFun φ) (p * u) G.measure ≤
        atomMeasureScale G s p uConj Q
\end{lstlisting}

\section{Besov-ish spaces}\label{besovish}

\filebox{WeakGrid/BesovishSpaces.lean}{WeakGrid/Atoms.lean}

This is the core file of the library. The Besov-ish space is defined by
\emph{declaring} what an atomic decomposition is and taking the functions
that admit one: a representation is a choice, for every level $k$ and cell
$Q$, of a coefficient $s_Q$ and an atom $a_Q$, such that the levelwise sums
converge to $g$ in $L^p$; its \emph{cost} is the $\ell^q$ norm over levels
of the $\ell^p$ norms over the cells of each level --- precisely the
sequence-space norm of the classical Frazier--Jawerth atomic decomposition,
transplanted to an abstract grid. Since representations are far from
unique, the norm is the infimum of the cost over all of them. Two analytic
inputs drive the section. The bounded overlap $G_1$ gives the levelwise
estimate (6-5) of the paper, and the summability hypothesis $G_2$ (a
weighted geometric-type series over the levels) lets the levelwise
estimates be summed; together they yield the embedding
$\mathcal{B}^s_{p,q}(\mathcal{A}) \subset L^t$ of Proposition~\ref{lp},
which in particular ($t=p$) shows that a finite-cost series really does
define an element of $L^p$ and that the infimum norm separates points.
Convexity and phase invariance of the atom sets then give the triangle
inequality, so the cost gauge is a norm.

Let $p \in [1,\infty)$, $u \in [1,\infty]$, $q \in [1,\infty]$, $s>0$,
$\mathcal{P}=(\mathcal{P}^k)_{k\geq 0}$ be a grid and let $\mathcal{A}$ be a
family of atoms of type $(s,p,u)$. We will also assume
\begin{itemize}
\item[$G_2$.] We have
$$\Big( \sum_k |\mathcal{P}^k|^{s q^\ast}\Big)^{1/q^\ast} < \infty
\quad\text{(usual modification if } q^\ast=\infty\text{)},
\qquad \lim_k |\mathcal{P}^k|=0,$$
where $q^\ast$ is the exponent produced by the H\"older-like estimate of the
paper (for $q>1$, $q^\ast = q'$).
\end{itemize}

\formalpar{$G_2$ is the predicate \lean{AssumptionG2} in
\lean{WeakGrid/BesovishSpaces.lean}, combining the finiteness of the
coefficient-weight series needed for the $t=p$ embedding with
$\lim_k|\mathcal{P}^k| = 0$. The paper writes this condition as
$C(p,q,(|\mathcal{P}^k|^s)_k)<\infty$ using the constant of its
Proposition~4.1; since paper Section~4 has no standalone Lean counterpart,
the formalization states the required summability of the weighted level
measures directly.}

The {\bf Besov-ish space} $\mathcal{B}^s_{p,q}(I,\mathcal{P},\mathcal{A})$ is
the set of all complex valued functions $g \in L^p$ that can be represented by
an absolutely convergent series in $L^p$
\begin{equation} \label{rep} g = \sum_{k=0}^{\infty} \sum_{Q \in \mathcal{P}^k} s_Q a_Q\end{equation}
where $a_Q$ is in $\mathcal{A}(Q)$, $s_Q \in \mathbb{C}$ and with finite
{\bf cost}
\begin{equation} \label{rep2} \big( \sum_{k=0}^{\infty} (\sum_{Q \in \mathcal{P}^k} |s_Q|^p)^{q/p} \big)^{1/q} < \infty\end{equation}
(usual modification when $q = \infty$). The series in (\ref{rep}) is called a
{\bf $\mathcal{B}^s_{p,q}(I,\mathcal{P},\mathcal{A})$-representation} of the
function $g$. Define
$$|g|_{\mathcal{B}^s_{p,q}(I,\mathcal{P},\mathcal{A})} = \inf \big( \sum_{k=0}^{\infty} (\sum_{Q \in \mathcal{P}^k} |s_Q|^p )^{q/p} \big)^{1/q},$$
where the infimum runs over all possible representations of $g$ as in
(\ref{rep}).

\formalpar{All in \lean{WeakGrid/BesovishSpaces.lean}. For each level $k$,
the structure \lean{LevelBlock} chooses one coefficient and one atom per
cell of $\mathcal{P}^k$; its value in $L^p$ (the finite sum
$\sum_{Q\in\mathcal{P}^k} s_Q a_Q$) is \lean{LevelBlock.toLp}. A
representation (\ref{rep}) is the structure \lean{LpGridRepresentation}: a
sequence of level blocks together with a \lean{HasSum} witness in $L^p$.
The cost (\ref{rep2}) is \lean{pqCost} (with the $\mathbb{R}_{\geq 0}\cup\{\infty\}$-valued
variant \lean{pqCostENNReal}), and its finiteness is the predicate
\lean{FinitePQCost}, with the endpoint $q=\infty$ represented by a supremum
bound. Membership with some finite-cost representation is
\lean{MemBesovishCoeffCost}; \lean{BesovishSpace} is the corresponding
\lean{Submodule} of $L^p$; and the norm
$|\cdot|_{\mathcal{B}^s_{p,q}(I,\mathcal{P},\mathcal{A})}$ is the cost gauge
\lean{Norm\_Costpq}, an infimum over representation costs. Implementation
notes: convergence of the series (\ref{rep}) is expressed by mathlib's
\lean{HasSum} over the level index, inside the Banach space
\lean{Lp} (over $\mathbb{C}$, exponent $p$), which encapsulates the paper's ``absolutely convergent in
$L^p$'' once the levelwise sums are grouped into blocks; the hypotheses
$1 \leq p$, $1 \leq q$ are carried as \lean{Fact} instances so that
mathlib's $L^p$ machinery picks them up by instance resolution; and the
normed structure on the space is provided as a \emph{local} instance
(\lean{costNormedAddCommGroup}) rather than a global one, because the same
underlying $L^p$ element may belong to several Besov-ish spaces with
different gauges.}

When the choice of $I$ and/or $\mathcal{P}$ is obvious we write
$\mathcal{B}^s_{p,q}(\mathcal{P},\mathcal{A})$ or
$\mathcal{B}^s_{p,q}(\mathcal{A})$; whenever we write just
$\mathcal{B}^s_{p,q}$ we choose the Souza's atoms $\mathcal{A}^{sz}_{s,p}$ of
Section~\ref{secatom}.
\begin{lstlisting}[style=lean]
structure LpGridRepresentation
    (A : AtomFamily G s p u) (g : Lp ℂ p G.measure) where
  block : (k : ℕ) → LevelBlock A k
  hasSum : HasSum (fun k => (block k).toLp A) g

def BesovishSpace (A : AtomFamily G s p u) (q : ℝ≥0∞)
    [Fact (1 ≤ q)]
    : Submodule ℂ (Lp ℂ p G.measure) where
  -- Carrier: all `L^p` elements admitting a Besov-ish atomic representation with finite p q cost.
  carrier := { g | MemBesovishCoeffCost A q g }
  ...

noncomputable def Norm_Costpq
    (A : AtomFamily G s p u) (q : ℝ≥0∞) [Fact (1 ≤ q)]
    (g : BesovishSpace A q) : ℝ :=
  pqPseudoNorm A q g
\end{lstlisting}

\begin{proposition}[Embedding into $L^t$]\label{lp}
Assume $G_1$--$G_2$ and $A_1$--$A_4$. Let $t \in [p,\infty)$ be such that
\begin{equation}\label{c11} s-\frac{1}{p}+\frac{1}{t} \geq 0, \qquad p \leq t \leq pu,\end{equation}
and suppose
\begin{equation}\label{decay} C_{2}(t)= C_1^{1+1/t}\, C\big(t,q,(|\mathcal{P}^k|^{t(s-\frac1p+\frac1t)})_k\big) < \infty.\end{equation}
Then for every $g$ with a $\mathcal{B}^s_{p,q}(\mathcal{A})$-representation
(\ref{rep}) of cost $K$ the series converges in $L^t$ and
\begin{equation}\label{inclulp} |g|_{t} \leq C_2(t)\, K.\end{equation}
In particular
$$|g|_t \leq C_2(t)\, |g|_{\mathcal{B}^s_{p,q}(\mathcal{A})},$$
and $\mathcal{B}^s_{p,q}(\mathcal{A}) \subset L^t$
continuously.
\end{proposition}

\formalpar{The theorem \lean{lp\_embedding\_adapted\_statement} in \\
\lean{WeakGrid/BesovishSpaces.lean} proves the per-representation inequality
(\ref{inclulp}): for any finite-cost representation $R$ of $g$, the $L^t$
norm of $g$ is bounded by the explicit constant
$C_1^{1+1/t}\cdot$\lean{cCoefficient} times the cost of $R$. The norm
inequality follows by taking the infimum over representations. Compared with
paper Proposition~6.1 the hypotheses are identical except for $q \geq 1$ and
$t \geq p$ being imposed globally, and the levelwise estimate (6.2) of the
paper is internal to the proof rather than part of the statement.}
\begin{lstlisting}[style=lean]
theorem lp_embedding_adapted_statement
    {A : AtomFamily G s p u} {t : ℝ≥0∞}
    [Fact (1 ≤ t)]
    (hp_ne_top : p ≠ ∞) (ht_ne_top : t ≠ ∞)
    (hq_one : 1 ≤ q)
    (hp_le_t : p ≤ t) (ht_le_pu : t ≤ p * u)
    (hs_nonneg : 0 ≤ s - 1 / p.toReal + 1 / t.toReal)
    {g : Lp ℂ p G.measure} (R : LpGridRepresentation A g)
    (hRfin : LpGridRepresentation.FinitePQCost (q := q) R)
    (hCco_fin : cCoefficientFinite t q (fun k =>
      (levelMeasureWeight G s p t k) ^ t.toReal)) :
        (MeasureTheory.eLpNorm (g : α → ℂ) t G.measure).toReal ≤
        ((G.grid.Cmult1 : ℝ) ^ (1 + 1 / t.toReal)) *
        cCoefficient t q (fun k => (levelMeasureWeight G s p t k) ^ t.toReal) *
          LpGridRepresentation.pqCost (q := q) R := ...
\end{lstlisting}

\begin{proposition}[Linear structure]\label{linear}
Assume $G_1$--$G_2$ and $A_1$--$A_4$. Then
$\mathcal{B}^s_{p,q}(\mathcal{A})$ is a complex linear subspace of $L^p$ and
$|\cdot|_{\mathcal{B}^s_{p,q}(\mathcal{A})}$ is a norm on it. Moreover the
inclusion
$$\imath\colon (\mathcal{B}^s_{p,q}(\mathcal{A}),|\cdot|_{\mathcal{B}^s_{p,q}(\mathcal{A})})\rightarrow (L^p,|\cdot|_p)$$
is continuous.
\end{proposition}

\formalpar{\lean{BesovishSpace} is by construction a \lean{Submodule} of
$L^p$, which gives the linear structure. The normed structure is the
definition \\ \lean{BesovishSpace.costNormedAddCommGroup} (requiring
$p \neq \infty$ and $G_2$), which installs \lean{Norm\_Costpq} as a norm
--- nonnegativity, triangle inequality and homogeneity are packaged, together
with the $L^t$ embedding estimate, in the summary theorem
\lean{BesovishSpace.normedSpace\_and\_lp\_embedding}. Continuity of $\imath$
is the case $t=p$ of Proposition~\ref{lp}. Compared with paper
Proposition~6.3, $\rho=\min\{1,p,q\}=1$ since $p,q\geq 1$, so the
$\rho$-norm of the paper is a norm.}
\begin{lstlisting}[style=lean]
noncomputable def costNormedAddCommGroup
    (hp_top : p ≠ ∞)
    (hCco_fin : LpGridRepresentation.cCoefficientFinite p q (fun k =>
      (LpGridRepresentation.levelMeasureWeight G s p p k) ^ p.toReal)) :
    NormedAddCommGroup (BesovishSpace A q) where
  norm := Norm_Costpq A q
  dist x y := Norm_Costpq A q (-x + y)
  ...
\end{lstlisting}

\section{Limits of representations, compactness, and completeness}
\label{complsec}

\filebox{WeakGrid/Completeness.lean}{WeakGrid/Atoms.lean,
WeakGrid/BesovishSpaces.lean}

The compactness and completeness theory rests on a single mechanism, the
\emph{limit of representations}: if functions $g_n$ have representations
with uniformly bounded cost, one extracts --- by a Cantor diagonal argument
over the countably many cells --- a subsequence along which every
coefficient converges and, using the compactness of the atom sets ($A_5$ or
$A_6$), every atom converges as well; the limiting coefficients and atoms
then form a representation of the limit function, and lower semicontinuity
of the cost under cellwise convergence gives the same cost bound $C$. This
one proposition yields, in cascade: sequential compactness of cost balls in
the $L^p$ topology, completeness of the Besov-ish space (a Cauchy sequence
converges in $L^p$ by the embedding, and the limit stays in the space with
the right norm), and later all the good-grid compactness statements of
Part~II. A technical point that the formalization makes explicit: to even
\emph{state} the limit representation one must know that a block sequence
with finite abstract cost is summable in $L^p$; this is
\lean{formalBlockSeq\_summable}, proved from $G_2$, and it is the reason
$G_2$ appears in all the statements of this section.

\begin{proposition}[Limits of representations]\label{compa2}
Assume $G_1$--$G_2$ and $A_1$--$A_4$. Suppose that $g_n$ are functions in
$\mathcal{B}^s_{p,q}(\mathcal{A})$ with
$\mathcal{B}^s_{p,q}(\mathcal{A})$-representations
$$g_n=\sum_{k=0}^{\infty}\sum_{Q\in \SP^k}s_Q^na_Q^n$$
satisfying
\begin{itemize}
\item[i.] There is $C$ such that for every $n$
$$\big(\sum_{k=0}^{\infty}(\sum_{Q\in \SP^k}|s_Q^n|^p)^{q/p}\big)^{1/q}\leq C.$$
\item[ii.] For every $Q\in \mathcal{P}$ the limit $s_Q=\lim_n s_Q^n$ exists.
\item[iii.] For every $Q\in \mathcal{P}$ there is $a_Q\in \mathcal{A}(Q)$ such
that either $a_Q^n \to a_Q$ in the strong topology of $L^p$, or $a_Q^n \to
a_Q$ weakly in $L^p$.
\end{itemize}
Then $g_n$ converges (strongly, respectively weakly) in $L^p$ to a function
$g \in \mathcal{B}^s_{p,q}(\mathcal{A})$ which has the
$\mathcal{B}^s_{p,q}(\mathcal{A})$-representation
$$g=\sum_{k=0}^{\infty}\sum_{Q\in \SP^k}s_Q\, a_Q$$
of cost at most $C$.
\end{proposition}

\formalpar{In \lean{WeakGrid/Completeness.lean}, matching paper
Proposition~6.4. The theorem \lean{representation\_limit} proves the weak
version: under uniform cost bounds (i), cellwise convergence of the
coefficients (ii) and weak convergence of the atoms (iii), the limiting
blocks define a Besov-ish function of cost at most $C$ and the $g_n$
converge weakly. The theorem \lean{representation\_limit\_strong} is the
strong $L^p$ version: strong convergence of the atoms gives strong
convergence of the represented functions.}
\begin{lstlisting}[style=lean]
theorem representation_limit
    (A : AtomFamily G s p u)(hG2 : AssumptionG2 G s p u q)
     {gseq : ℕ → Lp ℂ p G.measure}
    {gLim : Lp ℂ p G.measure} {C : ℝ}
    (H : RepresentationLimitHypotheses A q gseq gLim C)
    (hp_ne_top : p ≠ ∞) (hs_pos : 0 < s) (hu_one : 1 ≤ u)
    [Fact (1 ≤ u)] (hC : 0 ≤ C) :
    MemBesovishCoeffCost A q gLim ∧
      LpGridRepresentation.FinitePQCost (q := q) H.Rlim ∧
      LpGridRepresentation.pqCost (q := q) H.Rlim ≤ C ∧
      Tendsto (fun n => toWeakSpace ℂ (Lp ℂ p G.measure) (gseq n)) atTop
        (𝓝 (toWeakSpace ℂ (Lp ℂ p G.measure) gLim)) := ...
\end{lstlisting}

\begin{corollary}[Compactness and completeness]\label{compa1}
Assume $G_1$--$G_2$, $A_1$--$A_4$, and either $A_5$ or $A_6$. Then
\begin{itemize}
\item[i.] If $g_n \in \mathcal{B}^s_{p,q}(\mathcal{A})$ satisfy
$|g_n|_{\mathcal{B}^s_{p,q}(\mathcal{A})}\leq C$ for every $n$, then there is
a subsequence converging (strongly under $A_5$, weakly under $A_6$) in $L^p$
to some $g \in \mathcal{B}^s_{p,q}(\mathcal{A})$ with
$|g|_{\mathcal{B}^s_{p,q}(\mathcal{A})}\leq C$.
\item[ii.] Under $A_5$, every closed ball of
$(\mathcal{B}^s_{p,q}(\mathcal{A}),|\cdot|_{\mathcal{B}^s_{p,q}(\mathcal{A})})$
is sequentially compact in the strong topology of
$L^p$.
\item[iii.] Under $A_5$,
$(\mathcal{B}^s_{p,q}(\mathcal{A}),|\cdot|_{\mathcal{B}^s_{p,q}(\mathcal{A})})$
is a complex Banach space.
\end{itemize}
\end{corollary}

\formalpar{All in \lean{WeakGrid/Completeness.lean}. Item i.\ is \\
\lean{representation\_limit\_strong\_existence} and \\
\lean{representation\_limit\_weak\_existence} (paper Corollary~6.5(i)): they
start from a bare sequence of limit blocks, prove it is summable, and
package the resulting limit as a genuine representation of cost at most
$C$. Item ii.\ is \\
\lean{exists\_strongly\_convergent\_subseq\_of\_uniform\_pqCost} (extraction
of a strongly convergent subsequence from any uniformly cost-bounded
sequence; the extended version also returns the convergence of the
coefficients and of the atoms) and \\
\lean{closed\_Norm\_Costpq\_ball\_strongly\_seqCompact} (sequential
compactness of closed cost balls in the strong $L^p$ topology); paper
Corollary~6.5(iii) states the compactness of the inclusion $\imath$ under
$A_5$, and the formalization states the sequential compactness of closed
cost balls, which is the equivalent form used downstream. Item iii.\ is
\lean{besovishSpace\_costNorm\_completeSpace}, completeness of the Besov-ish
space for the cost norm (paper Corollary~6.5(ii); for $p,q \geq 1$ the
$\rho$-Banach space of the paper is a Banach space, and paper
Corollary~6.7 is the finite-dimensional criterion through which the Souza
case of Part~II is obtained). The paper's Corollary~6.6 (existence of a
representation attaining the infimum cost) is not formalized and is
omitted.}
\begin{lstlisting}[style=lean]
theorem closed_Norm_Costpq_ball_strongly_seqCompact
    (hp_ne_top : p ≠ ∞) (hs_pos : 0 < s) (hu_one : 1 ≤ u)
    (A : AtomFamily G s p u) (hG2 : AssumptionG2 G s p u q)
    (hA5 : AssumptionA5 A)
    {C : ℝ}
    [Fact (1 ≤ u)]
    (hC : 0 ≤ C)
    (gseq : ℕ → BesovishSpace A q)
    (hball : ∀ n, BesovishSpace.Norm_Costpq A q (gseq n) ≤ C) :
    ∃ (φ : ℕ → ℕ) (_hφ : StrictMono φ) (gLim : BesovishSpace A q),
      BesovishSpace.Norm_Costpq A q gLim ≤ C ∧
        Tendsto
          (fun n => ((gseq (φ n) : BesovishSpace A q) : Lp ℂ p G.measure))
          atTop
          (𝓝 ((gLim : BesovishSpace A q) : Lp ℂ p G.measure)) := ...

theorem besovishSpace_costNorm_completeSpace
    (hp_ne_top : p ≠ ∞) (hs_pos : 0 < s) (hu_one : 1 ≤ u)
    (A : AtomFamily G s p u) (hG2 : AssumptionG2 G s p u q)
    (hA5 : AssumptionA5 A)
    [Fact (1 ≤ u)] :
    @CompleteSpace (BesovishSpace A q)
      (BesovishSpace.costNormedAddCommGroup
        (A := A) (q := q) hp_ne_top hG2.1).toMetricSpace.toPseudoMetricSpace.toUniformSpace := ...
\end{lstlisting}

\section{Scales of spaces}\label{escala}

\filebox{WeakGrid/Scales.lean}{WeakGrid/BesovishSpaces.lean}

Changing the smoothness parameter from $s$ to $\tilde s$ amounts to
rescaling every atom of a cell $Q$ by the deterministic factor
$|Q|^{\tilde s - s}$ --- there is no analysis in the rescaling itself, only
in comparing the resulting spaces. When $s \leq \tilde s$ and the level
measures decay fast enough (a weighted summability over levels, the same
shape as $G_2$), a representation in the finer scale converts levelwise
into a representation in the coarser scale with a summable loss, giving the
continuous inclusion
$\mathcal{B}^{\tilde s}_{p,\tilde q} \subset \mathcal{B}^{s}_{p,q}$ ---
the abstract analogue of the classical monotonicity of Besov scales in the
smoothness parameter.

A family of atoms $\mathcal{A}$ of type $(s,p,u)$ induces a one-parameter
scale $\tilde{s}\mapsto \mathcal{A}_{\tilde{s},p}$, where
$\mathcal{A}_{\tilde{s},p}$ is the family of atoms of type $(\tilde{s},p,u)$
defined by
$$\mathcal{A}_{\tilde{s},p}(Q)= \{ |Q|^{\tilde{s}-s}a_Q\colon \ a_Q\in \mathcal{A}(Q)\}.$$

\begin{proposition}[Scale inclusion]\label{compa}
Assume $G_1$--$G_2$ and that $\mathcal{A}$ satisfies $A_1$--$A_4$. Let
$0< s \leq \tilde{s}$ and $q,\tilde{q} \in [1,\infty]$. Suppose
$$\Big( \sum_{k} |\mathcal{P}^k|^{q^\ast(\tilde{s}-s)} \Big)^{1/q^\ast} < \infty.$$
Then $\mathcal{B}^{\tilde{s}}_{p,\tilde{q}}(\mathcal{A}_{\tilde{s},p})\subset
\SB^s_{p,q}(\mathcal{A}_{s,p})$, the inclusion is a continuous linear
map, and
$$|g|_{\SB^s_{p,q}(\mathcal{A}_{s,p})} \leq
C\big(q,(|\mathcal{P}^k|^{\tilde s - s})_k\big)\,
|g|_{\mathcal{B}^{\tilde{s}}_{p,\tilde{q}}(\mathcal{A}_{\tilde{s},p})}.$$
\end{proposition}

\formalpar{In \lean{WeakGrid/Scales.lean}. The cellwise rescaling by the
deterministic factor $\mu(Q)^{\tilde s - s}$ is implemented by \\
\lean{smoothnessScaleFactor}, \lean{smoothnessScaleAtoms} and
\lean{smoothnessScaleAtomFamily}. The theorem \\
\lean{smoothnessScaleBesovishSpace\_subset} proves the set inclusion under
the deterministic summability hypothesis \lean{scaleCoefficientFinite} and
$s \leq \tilde s$; the definition
\lean{smoothnessScaleBesovishSpaceInclusion} packages it as a linear map,
and the theorem
\lean{smoothnessScaleBesovishSpaceInclusion\_Norm\_Costpq\_le} gives the
norm bound with the explicit constant \lean{scaleCoefficient}. This is
item~A of paper Proposition~7.1; items~B (compactness of bounded sets of
the finer space in the coarser norm) and C (compactness of the inclusion
under $A_7$) are not formalized and are omitted. For $u=\infty$ the
two-parameter scaling map \\
\lean{smoothnessIntegrabilityScaleAtomFamily} (scale factor
$\mu(Q)^{\tilde s - s + 1/p - 1/\tilde p}$) is also defined, but the paper's
two-parameter scale inclusion is not formalized.}
\begin{lstlisting}[style=lean]
theorem smoothnessScaleBesovishSpace_subset
  [Fact (1 ≤ p)]
  [Fact (1 ≤ u)] [Fact (1 ≤ q)] [Fact (1 ≤ qTilde)]
    (A : AtomFamily G s p u) (sTilde : ℝ) (hsTilde_pos : 0 < sTilde)
    (hs_le : s ≤ sTilde)
    (hqfin : scaleCoefficientFinite q (smoothnessScaleLevelWeight G s sTilde p))
    {g : Lp ℂ p G.measure}
    (hg : g ∈ SmoothnessScaleBesovishSpace
      (A.smoothnessScaleAtomFamily sTilde hsTilde_pos) qTilde) :
    g ∈ BesovishSpace A q := ...
\end{lstlisting}

\section{Induced spaces}\label{inducedsec}

\filebox{WeakGrid/InducedGrid.lean}{WeakGrid/BesovishSpaces.lean}

Induced spaces formalize the idea of \emph{zooming into a cell}: fixing a
cell $Q$ at level $k_0$, the cells of deeper levels contained in $Q$ form a
grid of the measure space $(Q, m|_Q)$, with levels re-indexed to start at
$0$, and the atoms supported in $Q$ form an atom family on this induced
grid. Reading a representation on the induced grid as a representation on
the ambient grid costs nothing (the coefficients are the same, placed at
levels $k_0 + i$), which is why the inclusion below is a weak contraction.
This construction is the basic localization tool of the library: Besov's
atoms (Section~\ref{alt2}) are \emph{defined} through induced norms, and
the multiplier theory (Section~\ref{multsec}) constantly moves between a
cell and the ambient space through it. The paper introduces induced spaces
for good grids (Section~10 of the published version); the formalization
develops them for arbitrary weak grids, since no good-grid geometry is
needed.

Consider a Besov-ish space $\mathcal{B}^s_{p,q}(I,\mathcal{P},\mathcal{A})$,
where $\mathcal{P}$ is a grid. Given $Q\in \mathcal{P}^{k_0}$, consider
the sequence of finite families of subsets
$\mathcal{P}_Q=(\mathcal{P}_Q^i)_{i\geq 0}$ of $Q$ given by
$$\mathcal{P}^i_Q=\{ P\in \mathcal{P}^{k_0+i},\ P\subset Q\},$$
and let $\mathcal{A}_Q$ be the restriction of $\mathcal{A}$ to indices in
$\mathcal{P}_Q$. Then we can consider the {\bf induced} Besov-ish space
$\mathcal{B}^s_{p,q}(Q,\mathcal{P}_Q,\mathcal{A}_Q)$, and the inclusion
$$i\colon \mathcal{B}^s_{p,q}(Q,\mathcal{P}_Q,\mathcal{A}_Q) \rightarrow \mathcal{B}^s_{p,q}(I,\mathcal{P},\mathcal{A})$$
is well defined and a weak contraction:
$$|f|_{\mathcal{B}^s_{p,q}(I,\mathcal{P},\mathcal{A})}\leq
|f|_{\mathcal{B}^s_{p,q}(Q,\mathcal{P}_Q,\mathcal{A}_Q)}.$$

\formalpar{In \lean{WeakGrid/InducedGrid.lean}. The induced grid and its
ambient space are the definitions \lean{inducedWeakGrid} and
\lean{inducedWeakGridSpace}; the restricted atom family is
\lean{inducedAtomFamily}; and the transfer of a representation from the
induced grid to the ambient one, with control of cost (which gives the weak
contraction above), is \lean{inducedRepresentationToAmbient}. For Souza's
atoms the formalization also transfers positivity along this map \\ 
(\lean{inducedRepresentationToAmbient\_atom\_toFunction\_pos}); this is the
form needed by the multiplier theory of Section~\ref{multsec}. The reverse restriction map is treated, for cells
of the grid, in Lemma~\ref{incint3} below.}
\begin{lstlisting}[style=lean]
def inducedRepresentationToAmbient
    (G : WeakGridSpace (α := α)) {k₀ : ℕ} (Q : LevelCell G k₀)
    {s : ℝ} {p u : ℝ≥0∞} [Fact (1 ≤ p)] (A : AtomFamily G s p u)
    {f : Lp ℂ p (inducedWeakGridSpace G Q).measure}
    (R : LpGridRepresentation (inducedAtomFamily G Q A) f) :
    LpGridRepresentation A (inducedLpToAmbient G Q f) := ...
\end{lstlisting}

\section{Transmutation of atoms}\label{transsec}

\filebox{WeakGrid/Transmutation.lean}{WeakGrid/Atoms.lean,
WeakGrid/BesovishSpaces.lean, WeakGrid/Completeness.lean}

The key result in Part~I (paper Proposition~8.1) states that if we replace
each atom of an atomic decomposition by a combination of atoms of a
possibly different class, in such way that we do not change the location of
the supports and the mass of the new representation concentrates pretty
much on the original scale, then we obtain a new atomic decomposition whose
cost is controlled by the cost of the original one.

The mechanism is worth describing, because the formalization follows it
closely. Each source atom $a_Q$, $Q \in \mathcal{G}^i$, is expanded in
target atoms $b_{P,Q}$ with coefficients $s_{P,Q}$ whose levelwise
$\ell^p$ mass decays geometrically away from the level $k_i$ attached to
$i$ (hypothesis III); the level selector $i \mapsto k_i$ is only assumed
\emph{almost linear} (hypothesis II), which is the right invariance for
comparing two grids whose scales are exponentially related. Summing over
$Q$ with coefficients $c_Q$ and collecting, for each target cell $P$, all
contributions, one obtains target coefficients
$m_P = \sum_Q |c_Q s_{P,Q}|$ and normalized target atoms $d_P$ (a convex
combination, by $A_3$ again). Estimating the cost of $(m_P)_P$ is then a
discrete convolution inequality across levels: the geometric decay
$\lambda^{k-k_i}$ acts as a convolution kernel in the level variable, and
the paper's convolution trick (Proposition~4.2) bounds the $\ell^{q/p}$
norm of the convolution by the kernel mass times the source cost. This is
where the explicit constant
$C_1 C_6^{1/p} \lambda^{-B/p} C_3(p,q,b) C_7^{1/q}$ of (8-17) comes from.
The truncated sums (over source levels $i \leq N$) are exact identities of
finite character (Claim~I); each truncation is a genuine representation
with the right cost (Claim~II); and the passage $N \to \infty$ uses the
limit-of-representations theorem of Section~\ref{complsec} (Claim~III) ---
this is why the file imports \lean{Completeness}.

\begin{proposition}[Transmutation of atoms]\label{trans}
Assume
\begin{itemize}
\item[{\bf I.}] $\mathcal{A}_2$ is a class of $(s,p,u_2)$-atoms for a grid
$\mathcal{W}$, satisfying $A_1$--$A_4$ and $G_1$--$G_2$. Let $\mathcal{G}$ be
also a grid satisfying $G_1$--$G_2$.
\item[{\bf II.}] $k_i\in \mathbb{N}$, $i\geq 0$, is a sequence such that there
are $\alpha>0$ and $A,B\in \mathbb{R}$ with
$$ \alpha i +A \leq k_i\leq \alpha i +B \quad\text{for every } i.$$
\item[{\bf III.}] There is $\lambda \in(0,1)$ such that the following holds.
For every $Q\in \mathcal{G}^i$ there are atoms $b_{P,Q} \in \mathcal{A}_2(P)$
and coefficients $s_{P,Q}\in \mathbb{C}$, indexed by the cells
$P \in \mathcal{W}$ with $P\subset Q$, such that
$$h_Q=\sum_k \sum_{P\in \mathcal{W}^k,\, P\subset Q} s_{P,Q} b_{P,Q}$$
is a $\mathcal{B}^s_{p,q}(\mathcal{A}_2)$-representation of a function $h_Q$,
with $s_{P,Q}=0$ whenever $P \in \mathcal{W}^k$ with $k < k_i$, and moreover
\begin{equation}\label{ad} \sum_{P\in \mathcal{W}^k,\, P \subset Q} |s_{P,Q}|^p \leq C_3\, \lambda^{k-k_i}
\quad\text{for every } k\geq k_i.\end{equation}
\end{itemize}
Then
\begin{itemize}
\item[{\bf A.}] For all coefficients $(c_Q)_{Q\in \mathcal{G}}$ with
$$\big( \sum_i \big( \sum_{Q\in \mathcal{G}^i} |c_{Q}|^p \big)^{q/p}\big)^{1/q}<\infty,$$
the sequence $N\mapsto \sum_{i\leq N} \sum_{Q\in \mathcal{G}^i }c_Q h_Q$
converges in $L^p$ to a function in $\mathcal{B}^s_{p,q}(\mathcal{A}_2)$ that
has a $\mathcal{B}^s_{p,q}(\mathcal{A}_2)$-representation
$\sum_k \sum_{P}m_P d_P$ with $m_P \geq 0$ and
$$\big(\sum_k \big( \sum_{P \in \mathcal{W}^k} |m_{P}|^p \big)^{q/p}\big)^{1/q}
\leq C_4\, \big( \sum_i \big( \sum_{Q\in \mathcal{G}^i} |c_{Q}|^p \big)^{q/p}\big)^{1/q},$$
where $C_4=C_4(C_1,C_3,\lambda,\alpha,A,B,p,q)$ is
explicit.
\item[{\bf B.}] Suppose, in addition to the assumptions of A., that the
$s_{P,Q}$ are nonnegative real numbers and $b_{P,Q}>0$ on $P$ for all $P,Q$.
If $m_P\neq 0$ and $d_P \neq 0$ somewhere on $P$, then $P\subset Q$ for some
cell $Q \in \mathcal{G}$ with $c_Q \neq 0$ and $s_{P,Q}>0$, and the
corresponding $h_Q$ is a positive multiple of a positive function a.e.\ on
$P$. If moreover $c_Q \geq 0$ for every $Q$, then $m_P\neq 0$ implies
$d_P > 0$ a.e.\ on $P$.
\item[{\bf C.}] Let $\mathcal{A}_1$ be a class of $(s,p,u_1)$-atoms for the
grid $\mathcal{G}$ satisfying $A_1$--$A_4$. Suppose that for every atom
$a_Q\in \mathcal{A}_1(Q)$ we can find $s_{P,Q}$ and $b_{P,Q}$ as in III.\ such
that $h_Q=a_Q$. Then
$$\mathcal{B}^s_{p,q}(\mathcal{A}_1)\subset \mathcal{B}^s_{p,q}(\mathcal{A}_2)$$
and this inclusion is continuous, with
$|\phi|_{\mathcal{B}^s_{p,q}(\mathcal{A}_2)} \leq C_4\,
|\phi|_{\mathcal{B}^s_{p,q}(\mathcal{A}_1)}$ for every $\phi \in
\mathcal{B}^s_{p,q}(\mathcal{A}_1)$.
\end{itemize}
\end{proposition}

\formalpar{In \lean{WeakGrid/Transmutation.lean}. The transmuted objects
(target coefficients $m_{P,N}$, normalized target atoms $d_{P,N}$, target
level blocks, and their stable limits) are the definitions
\lean{TransmutationCoeff}, \lean{TransmutationAtom},
\lean{TransmutationBlock} and \lean{TransmutationCoeffLimit},
\lean{TransmutationAtomLimit}, \lean{TransmutationBlockLimit}; the
hypothesis III is the predicate \lean{RepresentationWsubGandALS}, with the
almost-linear level control II packaged as \lean{AlmostLinearSequence}. The
internal pipeline is formalized by the theorems \lean{ClaimI} (exact
bookkeeping identity: the transmuted representation at truncation $N$ sums
to the source partial expansion), \lean{ClaimII} (each finite transmutation
block sequence is a genuine representation of the source partial expansion,
with cost controlled by the source cost) and \lean{ClaimIII} (passage to the
stable limiting data, producing the $L^p$ limit and its representation),
together with their endpoint versions \lean{ClaimI\_top},
\lean{ClaimII\_top}, \lean{ClaimIII\_top} for $q=\infty$. Item~A is the
theorem \lean{Transmutation\_of\_Atoms\_Claim\_A}, with the explicit
constant of the paper (product of $C_1$, $C_6^{1/p}$, $\lambda^{-B/p}$, the
coefficient of the convolution trick --- paper Proposition~4.2, which in
Lean is the constant \lean{cCoefficientInt} of the kernel
\lean{transmutationKernelZ} --- and the integer-part factor
$\max\{\ell\in\mathbb N : \ell < \alpha\}+1$). Item~B is
\lean{Transmutation\_of\_Atoms\_Claim\_B} and
\lean{Transmutation\_of\_Atoms\_Claim\_B\_sharp} (support and positivity of
the transmuted data); compared with paper Proposition~8.1.B, the pointwise
positivity conclusions are stated almost everywhere on $P$ (with positivity
witnessed by positive real multiples), which is the natural formulation in
$L^p$. Item~C is \lean{Transmutation\_of\_Atoms\_Claim\_C\_explicit} and the
continuous-embedding form
\lean{Transmutation\_of\_Atoms\_continuous\_embedding\_explicit}. For the
identity level selector, the theorem
\lean{LpGridRepresentation.cCoefficientInt\_transmutationKernelZ\_zero\_one}
computes the integer transmutation coefficient of the Claim~C kernel as the
geometric sum $(1-\rho)^{-1}$, $\lambda = \rho^p$; this is the constant
simplification used in the Souza/Besov atom comparison of
Section~\ref{alt2}.}
\begin{lstlisting}[style=lean,basicstyle=\ttfamily\scriptsize]
theorem Transmutation_of_Atoms_Claim_A
    (G W : WeakGridSpace (α := α))
    (AW : AtomFamily W s p u)
    (k : ℕ → ℕ)
    (A_als B_als r_als : ℝ)
    (hr_als : 0 < r_als)
    (hk_bound : ∀ i : ℕ,
      (k i : NNReal) ≤ r_als * (i : NNReal) + B_als ∧
      r_als * (i : NNReal) + A_als ≤ (k i : NNReal))
    (lam : ℝ) (hlam_pos : 0 < lam) (hlam_lt : lam < 1)
    (C : ℝ) (hC : 0 ≤ C)
    (h : (i : ℕ) → LevelCell G i → Lp ℂ p W.measure)
    (R : (i : ℕ) → (Q : LevelCell G i) → LpGridRepresentation AW (h i Q))
    (hR : RepresentationWsubGandALS (p := p) (q := q) G W AW k
      ⟨A_als, B_als, r_als, hr_als, hk_bound⟩ lam hlam_pos hlam_lt C hC h R)
    (c : (i : ℕ) → LevelCell G i → ℂ)
    (hc : CoeffFinitePQCost (p := p) (q := q) G c)
    (hq_ne_top : q ≠ ∞)
    (hG2_W : AssumptionG2 W s p u q)
    (hp_ne_top : p ≠ ∞)
    (hs_pos : 0 < s) :
    ∃ gLim : Lp ℂ p W.measure,
      ∃ hsum :
        HasSum (fun j => (TransmutationBlockLimit G W AW h R c A_als r_als j).toLp AW) gLim,
      MemBesovishCoeffCost AW q gLim ∧
      LpGridRepresentation.FinitePQCost (q := q)
        ({ block := TransmutationBlockLimit G W AW h R c A_als r_als
           hasSum := hsum } : LpGridRepresentation AW gLim) ∧
      Tendsto (fun N => PartialSumLevels G W h c N) atTop (𝓝 gLim) ∧
      LpGridRepresentation.pqCost (q := q)
        ({ block := TransmutationBlockLimit G W AW h R c A_als r_als
           hasSum := hsum } : LpGridRepresentation AW gLim) ≤
        (G.grid.Cmult1 : ℝ) *
        C ^ (1 / p.toReal) *
        lam ^ (-(B_als : ℝ) / p.toReal) *
        LpGridRepresentation.cCoefficientInt p ∞
          (transmutationKernelZ lam A_als r_als) *
        (Nat.ceil (r_als : ℝ) : ℝ) ^ (1 / q.toReal) *
        CoeffPQCost (p := p) (q := q) G c := ...

theorem Transmutation_of_Atoms_Claim_B (G W : WeakGridSpace (α := α))
    (AW : AtomFamily W s p u)
    (k : ℕ → ℕ)
    (A_als B_als r_als : ℝ)
    (hr_als : 0 < r_als)
    (hk_bound : ∀ i : ℕ,
      (k i : NNReal) ≤ r_als * (i : NNReal) + B_als ∧
      r_als * (i : NNReal) + A_als ≤ (k i : NNReal))
    (lam : ℝ) (hlam_pos : 0 < lam) (hlam_lt : lam < 1)
    (C : ℝ) (hC : 0 ≤ C)
    (h : (i : ℕ) → LevelCell G i → Lp ℂ p W.measure)
    (R : (i : ℕ) → (Q : LevelCell G i) → LpGridRepresentation AW (h i Q))
    (hR : RepresentationWsubGandALS_pos (p := p) (q := q) G W AW k
      ⟨A_als, B_als, r_als, hr_als, hk_bound⟩ lam hlam_pos hlam_lt C hC h R)
    (c : (i : ℕ) → LevelCell G i → ℂ)
    (hc : CoeffFinitePQCost (p := p) (q := q) G c)
    (hG2_W : AssumptionG2 W s p u q)
    (hp_ne_top : p ≠ ∞)
    (hs_pos : 0 < s) :
    (∃ gLim : Lp ℂ p W.measure,
      HasSum (fun j => (TransmutationBlockLimit G W AW h R c A_als r_als j).toLp AW) gLim ∧
      MemBesovishCoeffCost AW q gLim ∧
      Tendsto (fun N => PartialSumLevels G W h c N) atTop (𝓝 gLim)) ∧
    (∀ j : ℕ, ∀ P : LevelCell W j,
      TransmutationCoeffLimit G W AW h R c A_als r_als P ≠ 0 →
      (∃ x, x ∈ P.1 →
        AW.toFunction (levelCellToWeakGridCell W j P)
          (TransmutationAtomLocalLimit G W AW h R c A_als r_als P) x ≠ 0) →
        ∃ i ∈ Finset.range (transmutationStabilizationIndex A_als r_als j),
          ∃ Q ∈ (G.grid.partitions i).attach,
          P.1 ⊆ Q.1 ∧ c i Q ≠ 0 ∧
            ∃ r : NNReal, 0 < r ∧ ((R i Q).block j).coeff P = (r : ℂ) ∧
              ∀ᵐ x ∂ W.measure.restrict P.1,
                ∃ r : NNReal, 0 < r ∧ h i Q x = (r : ℂ)) := ...

theorem Transmutation_of_Atoms_continuous_embedding_explicit
    (G : WeakGridSpace (α := α))
    (u1 u2 : ℝ≥0∞)
    [Fact (1 ≤ u2)]
    (AG1 : AtomFamily G s p u1)
    (AG2 : AtomFamily G s p u2)
    (lam : ℝ) (hlam_pos : 0 < lam) (hlam_lt : lam < 1)
    (C : ℝ) (hC : 0 ≤ C)
    (hAG1_to_AG2 : ∀ i : ℕ, ∀ Q : LevelCell G i,
      ∀ g : Lp ℂ p G.measure,
        (∃ φ : (AG1.localSpace (levelCellToWeakGridCell G i Q)).carrier,
          AG1.IsAtom (levelCellToWeakGridCell G i Q) φ ∧
            g = atomLp AG1 (levelCellToWeakGridCell G i Q) φ) →
          ∃ Rg : LpGridRepresentation AG2 g,
            CoeffFinitePQCost (p := p) (q := q) G
              (fun j S => (Rg.block j).coeff S) ∧
            (∀ j : ℕ, ∀ S : LevelCell G j,
              (¬ S.1 ⊆ Q.1 → (Rg.block j).coeff S = 0) ∧
              (j < i → (Rg.block j).coeff S = 0)) ∧
            ∀ j : ℕ, i ≤ j → Rg.levelCoeffPower j ≤ C * lam ^ (j - i))
    (hG2_G : AssumptionG2 G s p u2 q)
    (hp_ne_top : p ≠ ∞)
    (hs_pos : 0 < s) :
    0 ≤ transmutationClaimCEmbeddingConstant G p q lam C ∧
      ∀ g : BesovishSpace AG1 q,
        ∃ hg2 : MemBesovishCoeffCost AG2 q (g : Lp ℂ p G.measure),
          BesovishSpace.Norm_Costpq AG2 q
              (⟨(g : Lp ℂ p G.measure), hg2⟩ : BesovishSpace AG2 q) ≤
            transmutationClaimCEmbeddingConstant G p q lam C *
              BesovishSpace.Norm_Costpq AG1 q g := ...
\end{lstlisting}

\section{Further weak-grid infrastructure}\label{infra}

Two more files complete the weak-grid layer; they contain no statement of
the paper but are imported by Part~II.

\filebox{WeakGrid/Multipliers.lean}{WeakGrid/Completeness.lean,
WeakGrid/InducedGrid.lean}

This file defines, at the weak-grid level of generality, the predicates on
which the multiplier theory of Section~\ref{multsec} is built: a measurable
$g$ \emph{is a pointwise multiplier} of a Besov-ish space
(\lean{IsPointwiseMultiplier}: multiplication by $g$ maps the space to
itself boundedly for the cost gauge) and $g$ \emph{is in the selfs class}
(\lean{PointwiseSelfsClass}: the products of $g$ with all atoms have
uniformly bounded Besov-ish norm). Keeping these definitions abstract means
that the implication ``multiplier $\Rightarrow$ selfs class'', which holds
for any atom family, lives here, while everything genuinely good-grid-ish
is deferred to \lean{GoodGrid/Multipliers/}.

\filebox{Sums.lean}{none}

Reusable indexing types and notation for sums over grid blocks: the global
pair type \lean{BlockIndex} for $(k,Q)$, the types \lean{ChildIndex} and
\lean{ParentIndex} for strict descendants and ancestors of a cell, and the
corresponding \lean{tsum} notations (\lean{sum\^{}b}, \lean{sum\^{}+},
\lean{sum\^{}-} in the source file).

\vspace{8mm}
\centerline{\Large\bf II. THE GOOD-GRID LIBRARY}
\smallskip
\centerline{\large\texttt{BesovSpacesGoodGrid/GoodGrid/}}
\vspace{6mm}

\begin{center}
\fbox{\begin{minipage}{0.88\textwidth}
\centering
In Part II we suppose that $s>0$, \ $p\in [1,\infty)$ \ and \
$q\in [1,\infty]$; from Section~\ref{multsec} on (the multiplier theory)
we further assume $0<s<1/p$.
\end{minipage}}
\end{center}
\vspace{4mm}

\section{Good grids}\label{goodgrids}

\filebox{GoodGrid/Definition.lean}{none (external:
\textnormal{\texttt{UnbalancedHaarWavelet}})}

A good grid adds to a weak grid the geometry of a nested dyadic-like
structure: levels partition the space ($G_4$--$G_5$), each cell sits inside
a parent at the previous level ($G_6$), and --- the quantitative heart ---
the measure of a child is comparable to the measure of its parent, with
ratio pinched between $\lambda_1$ and $\lambda_2$ ($G_7$). Two consequences
are used constantly. First, cell measures decay geometrically in the level,
$|Q| \leq \lambda_2^{k}|I|$ for $Q \in \mathcal{P}^k$, so the hypothesis
$G_2$ of Part~I holds \emph{automatically} for every $s > 0$: the entire
weak-grid theory applies to a good grid with no extra summability
assumptions. Second, a cell has at most $1/\lambda_1$ children, which
bounds all the branching constants in the Haar construction of
Section~\ref{haar}. The simplest example is the dyadic grid of $[0,1)$
with $\lambda_1 = \lambda_2 = 1/2$; a good grid, however, need not be
self-similar, and the two-sided ratio bound is the only regularity
retained.

A $(\lambda_1,\lambda_2)$-{\bf good grid}, with $0<\lambda_1\leq
\lambda_2<1$, is a grid $\mathcal{P}=(\mathcal{P}^k)_{k\in \mathbb{N}}$ with
the following properties:
\begin{itemize}
\item[$G_3$.] $\mathcal{P}^0=\{I\}$.
\item[$G_4$.] $I=\cup_{Q \in \mathcal{P}^k} Q$ up to a set of zero
$m$-measure.
\item[$G_5$.] The elements of $\mathcal{P}^k$ are pairwise disjoint.
\item[$G_6$.] For every $Q\in \mathcal{P}^k$, $k>0$, there exists
$P \in \mathcal{P}^{k-1}$ such that $Q\subset P$.
\item[$G_7$.] We have
$$\lambda_1 \leq \frac{|Q|}{|P|} \leq \lambda_2$$
for every $Q\subset P$ with $Q\in \mathcal{P}^{k+1}$ and
$P\in \mathcal{P}^{k}$, $k\geq 0$.
\item[$G_8$.] The family $\cup_k \mathcal{P}^k$ generates the
$\sigma$-algebra $\mathbb{A}$.
\end{itemize}

\formalpar{Good grids are the structure \lean{GoodGrid} in
\lean{GoodGrid/Definition.lean}, which extends the structure
\lean{UnbalancedHaarWavelet.Grid} of the external repository (where the
nested-partition properties $G_3$--$G_6$, $G_8$ live) with the ratio bounds
$G_7$ (fields \lean{ratio\_lower}, \lean{ratio\_upper}). The ambient object
used by the good-grid files is \lean{GoodGridSpace}, a measurable space with
a fixed good grid; the definition \lean{GoodGridSpace.toWeakGrid} turns a
good grid into a weak grid (since good-grid partitions are disjoint, the
same-level overlap multiplicity is $1$), so that the whole of Part~I
applies. The paper requires $\lambda_1 < \lambda_2$; the formalization
allows $\lambda_1 \le \lambda_2$.}
\begin{lstlisting}[style=lean]
structure GoodGrid extends UnbalancedHaarWavelet.Grid (α := α) where
        lambda1 : ℝ
        /-- Upper ratio bound (strictly less than 1). -/
        lambda2 : ℝ
        /-- `λ₁` is strictly positive. -/
        hlambda1_pos : 0 < lambda1
        /-- `λ₂` is strictly less than `1`. -/
        hlambda2_lt_one : lambda2 < 1
        /-- `λ₁ ≤ λ₂`. -/
        hlambda1_le_lambda2 : lambda1 ≤ lambda2
        /-- Each child cell's measure is at least `λ₁` times its parent's measure. -/
        ratio_lower : ∀ (n : ℕ) (s t : Set α),
                        s ∈ grid.partitions (n + 1) → t ∈ grid.partitions n → s ⊆ t →
                        ENNReal.ofReal lambda1 * μ t ≤ μ s
        /-- Each child cell's measure is at most `λ₂` times its parent's measure. -/
        ratio_upper : ∀ (n : ℕ) (s t : Set α),
                        s ∈ grid.partitions (n + 1) → t ∈ grid.partitions n → s ⊆ t →
                        μ s ≤ ENNReal.ofReal lambda2 * μ t
\end{lstlisting}

\section{Souza's atoms}\label{secatom}

\filebox{GoodGrid/BesovSpace.lean}{WeakGrid/Atoms.lean,
WeakGrid/BesovishSpaces.lean, WeakGrid/Completeness.lean,
WeakGrid/InducedGrid.lean, GoodGrid/Definition.lean}

Souza's atoms (paper Section~11A) are the simplest conceivable atoms: a
multiple of the indicator of a cell, normalized in $L^\infty$. Their local
model space is one-dimensional ($\mathcal{B}(Q) \cong \mathbb{C}$), so the
atom set of every cell is a closed bounded subset of a one-dimensional
space, hence compact: assumption $A_5$ comes for free (this is the
finite-dimensional criterion of paper Corollary~6.7), and the whole
compactness/completeness theory of Part~I applies. Their simplicity is the
point of the construction: products, restrictions, and positivity of
Souza's atoms are trivial to control, which is what makes the multiplier
theory of Section~\ref{multsec} tractable.

Let $Q\in \mathcal{P}$. A {\bf $(s,p)$-Souza's atom} supported on $Q$ is a
function $a\colon I \rightarrow \mathbb{C}$ such that $a(x)=0$ for every
$x \not\in Q$ and $a$ is constant on $Q$, with
$$|a|_\infty\leq |Q|^{s-1/p}.$$
The set of Souza's atoms supported on $Q$ will be denoted by
$\mathcal{A}^{sz}_{s,p}(Q)$. The {\bf canonical Souza's atom} on $Q$ is the
Souza's atom with $a(x)= |Q|^{s-1/p}$ for every $x \in Q$. Souza's atoms are
$(s,p,\infty)$-type atoms and satisfy $A_1$--$A_5$.

\formalpar{In \lean{GoodGrid/BesovSpace.lean}. The predicate
\lean{IsSouzaAtom} says that a function is supported on a good-grid cell,
constant on it, and of size at most $|Q|^{s-1/p}$; the canonical atom is
\lean{canonicalSouzaAtom}, with \lean{canonicalSouzaAtom\_isSouzaAtom}
proving that it is indeed a Souza's atom. The definitions
\lean{souzaLocalVectorSpace}, \lean{souzaAtomsSet} and
\lean{souzaAtomFamily} package Souza's atoms as an \lean{AtomFamily} with
$u=\infty$ satisfying $A_1$--$A_5$. The paper's further examples ---
H\"older atoms (Section~11B) and bounded variation atoms (Section~11C) ---
are not formalized and are omitted.}
\begin{lstlisting}[style=lean]
def IsSouzaAtom (G : GoodGridSpace (α := α)) (s : ℝ) (p : ℝ≥0∞)
    (Q : GoodGridCell G) (a : α → ℂ) : Prop :=
  (∀ x ∉ Q.cell, a x = 0) ∧
  ∃ c : ℂ, (∀ x ∈ Q.cell, a x = c) ∧
            ‖c‖ ≤ (G.grid.μ Q.cell).toReal ^ (s - (p.toReal)⁻¹)
\end{lstlisting}

\section{Besov spaces in a measure space with a good grid}\label{secbesov}

\filebox{GoodGrid/BesovSpace.lean}{(same file as the previous section)}

The paper's Section~12 introduces the spaces
$\mathcal{B}^s_{p,q} = \mathcal{B}^s_{p,q}(\mathcal{P},\mathcal{A}^{sz}_{s,p})$
and observes that the general theory applies; it states no numbered
theorem there. The theorem below is the formalization's packaging of the
resulting Souza-space facts, obtained by instantiating
Proposition~\ref{lp} and Corollary~\ref{compa1} (paper Proposition~6.1 and
Corollaries~6.5 and~6.7). The density statements in item~C deserve a
remark: a function constant on the cells of level $k$ is, by definition, a
finite sum of level-$k$ Souza's atoms, and as $k \to \infty$ such
``grid-simple'' functions approximate any $L^p$ function (this is where
$G_8$, the generation of the $\sigma$-algebra by the cells, enters); for
$s < 1/p$ the cost of these approximations stays finite, giving density
--- so, in contrast with classical Besov spaces of positive smoothness on
$\mathbb{R}^n$, here the space is a \emph{dense} compactly embedded
subspace of $L^p$.

We study the Besov-ish spaces
$\mathcal{B}^s_{p,q}=\mathcal{B}^s_{p,q}(\mathcal{P},\mathcal{A}^{sz}_{s,p})$
associated with a measure space with a good grid $(I,\mathcal{P},m)$. If
$p\in [1,\infty)$, $q\in [1,\infty]$ and $0<s<\frac{1}{p}$ we say that
$\mathcal{B}^s_{p,q}$ is a {\bf Besov space}. The general theory of Part~I
specializes as follows.

\begin{theorem}\label{souzamain}
Let $(I,\mathcal{P},m)$ be a measure space with a good grid, $s>0$,
$p \in [1,\infty)$, $q \in [1,\infty]$. Then
\begin{itemize}
\item[A.] $(\mathcal{B}^s_{p,q}, |\cdot|_{\mathcal{B}^s_{p,q}})$ is a complex
Banach space.
\item[B.] Closed balls of $\mathcal{B}^s_{p,q}$ are compact subsets of $L^p$,
and also of $L^1$.
\item[C.] If $0<s<1/p$ then $\mathcal{B}^s_{p,q}$ is dense in $L^p$, and
dense in $L^1$.
\end{itemize}
\end{theorem}

\formalpar{All in \lean{GoodGrid/BesovSpace.lean}. Item~A is
\lean{souzaBesovSpace\_costNorm\_completeSpace}, the specialization of the
weak-grid completeness theorem to Souza's atoms on a good grid. Item~B is
\lean{souza\_closedCostBallInLp\_isSeqCompact},
\lean{souza\_closedCostBallInLp\_isCompact} and
\lean{souza\_closedCostBallInL1\_isCompact}: compactness of closed cost
balls in the ambient $L^p$ topology and, after the finite-measure inclusion,
in $L^1$ --- the Souza-atom instance of Corollary~\ref{compa1} together with
the embedding of Proposition~\ref{lp}. Item~C is
\lean{souzaBesovSpace\_dense} and \lean{souzaBesovSpace\_dense\_inL1},
density of the Souza Besov space in $L^p$ and of its image in $L^1$; the
density statements do not appear explicitly as theorems in the paper, and
are recorded here because they are formalized.}
\begin{lstlisting}[style=lean,basicstyle=\ttfamily\scriptsize]
theorem souzaBesovSpace_costNorm_completeSpace
    (G : GoodGridSpace (α := α))
    (s : ℝ) (p q : ℝ≥0∞)
    (hs : 0 < s) (hp : 1 ≤ p) (hp_top : p ≠ ∞)
    [Fact (1 ≤ p)] [Fact (1 ≤ q)] :
    @CompleteSpace
      (SouzaBesovSpace G s p q hs hp hp_top)
      (WeakGridSpace.BesovishSpace.costNormedAddCommGroup
        (A := souzaAtomFamily G s p hs hp hp_top) (q := q)
        hp_top (souza_hCco G s p q hs hp hp_top)).toMetricSpace.toPseudoMetricSpace.toUniformSpace := ...

theorem souza_closedCostBallInLp_isCompact
    (G : GoodGridSpace (α := α))
    (s : ℝ) (p q : ℝ≥0∞)
    (hs : 0 < s) (hp : 1 ≤ p) (hp_top : p ≠ ∞)
    [Fact (1 ≤ p)] [Fact (1 ≤ q)]
    {C : ℝ} (hC : 0 ≤ C) :
    IsCompact
      {f : MeasureTheory.Lp ℂ p G.toWeakGridSpace.measure |
        ∃ hf : f ∈ SouzaBesovSpace G s p q hs hp hp_top,
          WeakGridSpace.BesovishSpace.Norm_Costpq
            (souzaAtomFamily G s p hs hp hp_top) q ⟨f, hf⟩ ≤ C} := ...

theorem souzaBesovSpace_dense
    (G : GoodGridSpace (α := α))
    (s : ℝ) (p q : ℝ≥0∞)
    (hs : 0 < s) (hp : 1 ≤ p) (hp_top : p ≠ ∞)
    [Fact (1 ≤ p)] [Fact (1 ≤ q)] :
    Dense (SouzaBesovSpace G s p q hs hp hp_top :
      Set (MeasureTheory.Lp ℂ p G.toWeakGridSpace.measure)) := ...
\end{lstlisting}

\section{Besov's atoms and the Souza--Besov comparison}\label{alt2}

\filebox{GoodGrid/BesovAtoms.lean}{GoodGrid/BesovSpace.lean,
WeakGrid/InducedGrid.lean, WeakGrid/Transmutation.lean}

This is the paper's ``Alternative characterizations, II: messing with
atoms'' (Section~16 of the published version), and the first serious
application of the transmutation machinery. A Besov atom on $Q$ is not
required to be constant --- it is any function supported on $Q$ whose
\emph{induced} Besov norm of a higher smoothness $\beta > s$ is controlled
by the scale $|Q|^{s-\beta}$; the induced spaces of
Section~\ref{inducedsec} are exactly what makes this definition possible.
The comparison theorem (paper Proposition~16.1) says that any atom family
sandwiched between Souza's atoms and Besov's atoms generates the
\emph{same} space with equivalent norms. The nontrivial inclusion is
proved by transmutation: each Besov atom is expanded, inside its own cell,
into Souza's atoms with geometrically decaying levelwise mass (rate
$\lambda_2^{(k-j)(\beta-s)}$ --- this is where $\beta > s$ and the
good-grid ratio bound enter), which is precisely hypothesis~III of paper
Proposition~8.1 with the identity level selector; Claim~C of that
proposition then converts every representation by Besov atoms into one by
Souza's atoms with controlled cost. In practice this proposition is the
standard tool for checking that a concrete function belongs to
$\mathcal{B}^s_{p,q}$: it suffices to exhibit atoms in any sandwiched
class.

Let $s < \beta$ and $\tilde{q}\in [1,\infty]$. A
{\bf $(s,\beta,p,\tilde{q})$-Besov atom} on $Q$ is a function
$a \in \mathcal{B}^\beta_{p,\tilde{q}}(Q,\mathcal{P}_Q,\mathcal{A}^{sz}_{\beta,p})$
such that $a(x)=0$ for $x \not\in Q$ and
$$|a|_{\mathcal{B}^\beta_{p,\tilde{q}}(Q, \mathcal{P}_Q, \mathcal{A}^{sz}_{\beta,p})}\leq \frac{1}{C_7} |Q|^{s-\beta},$$
where $C_7 = C_7(C_1,\lambda_2,\beta,\tilde q,p) \geq 1$ is the normalizing
constant of the paper. The family of $(s,\beta,p,\tilde{q})$-Besov atoms
supported on $Q$ is denoted $\mathcal{A}^{bs}_{s,\beta,p,\tilde{q}}(Q)$, and
$\mathcal{A}^{sz}_{s,p}(Q)\subset \mathcal{A}^{bs}_{s,\beta,p,\tilde{q}}(Q)$;
Besov atoms are atoms of type $(s,p,1)$.

\begin{proposition}[Souza's atoms and Besov's atoms]\label{besova}
Let $\mathcal{P}$ be a good grid and $0<s<1/p$. Let $\mathcal{A}$ be a class
of $(s,p,u)$-atoms, with $u\geq 1$, such that for some $s< \beta$,
$\tilde{q}\in [1,\infty]$, and $C_8, C_9 \geq 0$ we have, for every
$Q\in \mathcal{P}$,
$$\frac{1}{C_8} \mathcal{A}^{sz}_{s,p}(Q) \subset \mathcal{A}(Q)\subset C_9\, \mathcal{A}^{bs}_{s,\beta,p,\tilde{q}}(Q).$$
Then
$$\mathcal{B}^s_{p,q}(\mathcal{P},\mathcal{A}^{sz}_{s,p})=\mathcal{B}^s_{p,q}(\mathcal{P},\mathcal{A})=\mathcal{B}^s_{p,q}(\mathcal{P},\mathcal{A}^{bs}_{s,\beta,p,\tilde{q}}),$$
with
$$|f|_{\mathcal{B}^s_{p,q}(\mathcal{A})} \leq C_8\, |f|_{\mathcal{B}^s_{p,q}(\mathcal{A}^{sz}_{s,p})}
\quad\text{and}\quad
|f|_{\mathcal{B}^s_{p,q}(\mathcal{A}^{sz}_{s,p})} \leq C\, |f|_{\mathcal{B}^s_{p,q}(\mathcal{A})}.$$
\end{proposition}

\formalpar{This is paper Proposition~16.1. The sandwich proposition is the
theorem \lean{atoms\_between\_souza\_atoms\_and\_besov\_atoms} in
\lean{GoodGrid/BesovAtoms.lean}: it proves the equality of the three
Besov-ish spaces and both continuous norm bounds, the second with the
explicit constant $C_2/(1-\lambda_2^{\beta-s})$. It is supported by the
decay claims \lean{besovAtom\_to\_souza\_representation\_decay} and
\lean{besovAtom\_to\_induced\_souzaS\_representation\_decay\_claimC} (the
Lean counterparts of the Claim in the paper's proof: a single Besov atom,
read on the weak grid induced on its supporting cell, admits a Souza-atom
representation with geometric coefficient decay rate
$\lambda_2^{(k-j)(\beta-s)}$), by the normalizing constant
\lean{besovAtomConstant} (with \lean{one\_le\_besovAtomConstant} giving the
lower bound $1 \leq C_7$ used to simplify the final constant), and by the
identity-level geometric-sum computation
\lean{LpGridRepresentation.cCoefficientInt\_transmutationKernelZ\_zero\_one}
of Section~\ref{transsec}. Paper Propositions~16.2 (H\"older atoms) and
16.3 (bounded variation atoms) are not formalized and are omitted.}
\begin{lstlisting}[style=lean,basicstyle=\ttfamily\scriptsize]
theorem atoms_between_souza_atoms_and_besov_atoms
    (G : GoodGridSpace (α := α)) (s β : ℝ) (p u q qtilde : ℝ≥0∞)
    [Fact (1 ≤ p)] [Fact (1 ≤ u)] [Fact (1 ≤ q)] [Fact (1 ≤ qtilde)]
    (hs : 0 < s) (hβ : 0 < β) (hβs : s < β) (hp_top : p ≠ ∞)
    (A : WeakGridSpace.AtomFamily G.toWeakGridSpace s p u)
    (C1 C2 : ℝ)
    (hSandwich :
      SouzaBesovSandwich G s β p u qtilde hs hβ (inferInstance : Fact (1 ≤ p))
        hp_top A C1 C2) :
    (WeakGridSpace.BesovishSpace (souzaAtomFamily G s p hs (Fact.out : 1 ≤ p) hp_top) q =
        WeakGridSpace.BesovishSpace A q) ∧
      (WeakGridSpace.BesovishSpace A q =
        WeakGridSpace.BesovishSpace
          (besovAtomFamily G s β p qtilde hs hβ inferInstance hp_top) q) ∧
      (∀ f : WeakGridSpace.BesovishSpace
          (souzaAtomFamily G s p hs (Fact.out : 1 ≤ p) hp_top) q,
        ∃ hfA : WeakGridSpace.MemBesovishCoeffCost A q
            (f : Lp ℂ p G.toWeakGridSpace.measure),
          WeakGridSpace.BesovishSpace.Norm_Costpq A q
              (⟨(f : Lp ℂ p G.toWeakGridSpace.measure), hfA⟩ :
                WeakGridSpace.BesovishSpace A q)
            ≤ C1 *
              WeakGridSpace.BesovishSpace.Norm_Costpq
                (souzaAtomFamily G s p hs (Fact.out : 1 ≤ p) hp_top) q f) ∧
      (∀ f : WeakGridSpace.BesovishSpace A q,
        ∃ hfS : WeakGridSpace.MemBesovishCoeffCost
            (souzaAtomFamily G s p hs (Fact.out : 1 ≤ p) hp_top) q
            (f : Lp ℂ p G.toWeakGridSpace.measure),
          WeakGridSpace.BesovishSpace.Norm_Costpq
              (souzaAtomFamily G s p hs (Fact.out : 1 ≤ p) hp_top) q
              (⟨(f : Lp ℂ p G.toWeakGridSpace.measure), hfS⟩ :
                WeakGridSpace.BesovishSpace
                  (souzaAtomFamily G s p hs (Fact.out : 1 ≤ p) hp_top) q)
            ≤ (C2 / (1 - G.grid.lambda2 ^ (β - s))) *
              WeakGridSpace.BesovishSpace.Norm_Costpq A q f) := ...
\end{lstlisting}

\section{Unbalanced Haar wavelets}\label{haar}

\filebox{GoodGrid/AlternativeRepresentationsAndNorms/HaarRepresentationNorm.lean}{GoodGrid/BesovSpace.lean
(external: \textnormal{\texttt{UnbalancedHaarWavelet}},
\textnormal{\texttt{UnconditionalSchauderBasis}},
\textnormal{\texttt{Burkholder}})}

On a good grid there is a genuine wavelet system: since each cell has
between $2$ and $1/\lambda_1$ children, one can follow Girardi and
Sweldens and organize the children of each cell into a balanced binary
tree of subsets (the ``logarithmic subtrees''), attaching to each internal
node a two-valued zero-mean function --- an \emph{unbalanced Haar
function}. The measure-ratio bounds $G_7$ guarantee that each
$\phi_S$ supported in $Q$ has size $\approx |Q|^{-1/2}$ on its support, so
the system behaves like the classical Haar basis even though the grid may
be very irregular. The unconditionality of this basis in $L^\beta$,
$1 < \beta < \infty$ --- a Burkholder-type martingale inequality --- is
the one substantial analytic ingredient that the repository does not prove
internally: it is supplied by the external formalized libraries listed
above.

Let $\mathcal{P}=(\mathcal{P}^k)_k$ be a good grid. Following Girardi and
Sweldens, one constructs, from the children structure of the grid (Type I
tree with the logarithmic subtrees construction), a family of {\bf unbalanced
Haar functions} $\{\phi_S\}_{S\in \hat{\mathcal{H}}}$, where
$\hat{\mathcal{H}}=\{I\}\cup \mathcal{H}$, $\phi_I = 1_I/|I|^{1/2}$, each
$\phi_S$ with $S \in \mathcal{H}_Q$ is supported on $Q$, has zero mean,
$L^2$-norm one, and satisfies
$$\frac{C_5(\lambda_1,\lambda_2)}{|Q|^{1/2}} \leq |\phi_S(x)| \leq \frac{C_6(\lambda_1,\lambda_2)}{|Q|^{1/2}}$$
on its support. The family $\{\phi_S\}_{S\in \hat{\mathcal{H}}}$ is an
unconditional basis of $L^\beta$ for every $1<\beta<
\infty$.

\formalpar{The construction and the unconditional basis property are
provided by the external Lean repositories \lean{UnbalancedHaarWavelet},
\lean{UnconditionalSchauderBasis}, \lean{LaminarFamiliesMaximalBinaryTrees}
and \lean{Burkholder} (SmaniaD), on which this repository depends. The local
API ($L^2$-normalized Haar functions \lean{L2normalizedHaar} and Haar
coefficients \lean{Coeff}) is in
\lean{GoodGrid/AlternativeRepresentationsAndNorms/HaarRepresentationNorm.lean}.}
\begin{lstlisting}[style=lean]
def L2normalizedHaar (G : GoodGridSpace (α := α)) [DecidableEq (Set α)]
    (F : UnbalancedHaarWavelet.FullHaarSystem (G := GridOf G))
    (i : F.Index) (x : α) : ℂ :=
  ((l2NormalizationFactor G F i : ℝ) : ℂ) *
    (UnbalancedHaarWavelet.FullHaarSystem.function (GridOf G) F i x : ℂ)

def Coeff (G : GoodGridSpace (α := α)) [DecidableEq (Set α)]
    (F : UnbalancedHaarWavelet.FullHaarSystem (G := GridOf G))
    (f : α → ℂ) (_hf : Integrable f G.grid.μ) (i : F.Index) : ℂ :=
  ∫ x, f x * L2normalizedHaar G F i x ∂G.grid.μ
\end{lstlisting}

\section{Alternative norms: standard, Haar, and mean oscillation}\label{alt1}

\filebox{GoodGrid/AlternativeRepresentationsAndNorms/ (eleven files, see
Section~\ref{guide})}{GoodGrid/BesovSpace.lean, WeakGrid/InducedGrid.lean,
HaarRepresentationNorm.lean}

This is the paper's ``Alternative characterizations, I: messing with
norms'' (Section~15 of the published version). The Besov norm is an
infimum over all representations, which makes it hard to compute for a
concrete $f$; this section provides three norms defined directly from $f$
and proves them equivalent to it. The \emph{Haar norm} $N_{haar}$ weighs
the Haar coefficients $d_S^f = \int f\phi_S\,dm$; the \emph{standard
norm} $N_{st}$ weighs the coefficients of one distinguished atomic
representation, obtained by resumming the Haar expansion cellwise (each
$k_P^f$ is a finite, geometry-bounded combination of the $d_S^f$ with $S$
in the Haar tree of the parent of $P$); and the \emph{mean-oscillation
norm} $N_{osc}$ measures, in the spirit of Dorronsoro, how well $f$ is
approximated by constants at every cell and scale. The equivalences are
proved as a cycle
$|\cdot|_{\mathcal{B}^s_{p,q}} \to N_{st} \to N_{haar} \to N_{osc} \to
|\cdot|_{\mathcal{B}^s_{p,q}}$ (inequalities (15-30)--(15-33) of the
paper), each step with its own mechanism: testing $f - c_Q$ against the
zero-mean wavelets bounds Haar coefficients by oscillations; the constancy
of low-frequency Souza blocks on finer cells means oscillation only sees
the high-frequency tail of a representation, which a discrete convolution
estimate sums; and the standard representation, being an actual
representation, dominates the infimum norm. In particular the standard
representation converges to $f$ for \emph{every} $f$ in the space --- a
canonical, linear-in-$f$ decomposition that substitutes for the
nonuniqueness of general representations.

For $f\in L^\beta$, $\beta>1$, write $d_S^f = \int f \phi_S\,dm$ for the Haar
coefficients of $f$, and define
$$ N_{haar}(f)=|I|^{1/p-s-1/2} |d_I^f|+ \Big( \sum_{k} \big( \sum_{Q\in \mathcal{P}^{k}} |Q|^{1-sp-\frac{p}{2}} \sum_{S\in \mathcal{H}_Q} |d_S^f|^p\big)^{q/p} \Big)^{1/q}.$$
From the Haar coefficients one builds, exactly as in the paper, the
{\bf standard atomic representation}
$$f = \sum_i \sum_{Q\in \mathcal{P}^{i}} k_Q^f\, a_Q$$ ($a_Q$ canonical Souza's
atoms) and its norm
$$N_{st}(f)=|k_I^f|+ \Big( \sum_{k\geq 1} \big( \sum_{Q\in \mathcal{P}^{k}} |k_Q^f|^p \big)^{q/p}\Big)^{1/q}.$$
Finally, with
$$osc_p(f,Q)= \inf_{c \in \mathbb{C}} \Big(\int_Q |f-c|^p \, dm\Big)^{1/p},$$
define
$$N_{osc}(f)=|I|^{-s}|f|_p+\Big( \sum_k \big( \sum_{Q\in \mathcal{P}^k} |Q|^{-sp} osc_p(f,Q)^p \big)^{q/p} \Big)^{1/q}.$$

\begin{theorem}[Equivalence of norms]\label{alte}
Suppose $0<s<1/p$, $p\in [1,\infty)$ and $q\in [1,\infty]$. Each one of the
quantities $|f|_{\mathcal{B}^s_{p,q}}$, $N_{st}(f)$, $N_{haar}(f)$,
$N_{osc}(f)$ is finite if and only if $f\in \mathcal{B}^s_{p,q}$, and these
norms are equivalent on $\mathcal{B}^s_{p,q}$. More precisely there are
constants, depending only on the grid geometry and on $s,p,q$, such that
\begin{align*}
N_{st}(f) &\leq C\, |f|_{\mathcal{B}^s_{p,q}}, &
N_{st}(f) &\leq C\, N_{haar}(f),\\
N_{haar}(f) &\leq C\, N_{osc}(f), &
N_{osc}(f) &\leq C\, |f|_{\mathcal{B}^s_{p,q}},
\end{align*}
and conversely: if $N_{st}(f)<\infty$ then $f \in \mathcal{B}^s_{p,q}$ and
the standard atomic representation is a
$\mathcal{B}^s_{p,q}$-representation of $f$ (so
$|f|_{\mathcal{B}^s_{p,q}} \leq N_{st}(f)$); if $N_{haar}(f) < \infty$ and
$f \in L^1$ then $f \in L^p$ and the Haar series of $f$ converges in
$L^p$.
\end{theorem}

\formalpar{In the files of
\lean{GoodGrid/AlternativeRepresentationsAndNorms/}. The equivalences are
proved as a cycle of separate inequalities, each with a finite existential
constant:
\lean{exists\_standardRepresentationNorm\_le\_const\_mul\_souzaBesovNorm}
($N_{st} \leq C\,|\cdot|_{\mathcal{B}^s_{p,q}}$),
\lean{exists\_standardRepresentationNorm\_le\_const\_mul\_haarL2RepresentationNorm}
(the Haar gauge controls the standard atomic gauge),
\lean{exists\_haarL2RepresentationNorm\_le\_const\_mul\_meanOscillationNorm}
(mean oscillation controls the Haar gauge; the cellwise estimate is the Lean
analogue of testing $f - c_Q$ against the zero-mean Haar wavelets over a
parent cell) and
\lean{exists\_meanOscillationNorm\_le\_const\_mul\_souzaBesovNorm} (the
Besov norm controls mean oscillation; the proof separates the $L^p$ term
from the oscillation seminorm and uses that low-frequency Souza blocks are
constant on finer cells, so local oscillation is controlled by the
high-frequency tail and a geometric discrete convolution estimate). The
converse endpoints are
\lean{finite\_standardRepresentationNorm\_implies\_}\allowbreak\lean{memBesov\_and\_standardRepresentation}
(finite $N_{st}$ gives $L^p$ membership, the canonical standard
representation as a genuine finite-cost
$\mathcal{B}^s_{p,q}$-representation, and Souza-Besov membership with the
expected cost bounds) and
\lean{finite\_haarL2RepresentationNorm\_implies\_memLp\_and\_hasSum}
(finite $N_{haar}$ gives $f \in L^p$ and convergence of the normalized Haar
expansion to $f$ in $L^p$); the wrap-up of the Haar-side comparisons is
\lean{exists\_haarRepresentationNorms\_equivalent}. Compared with paper
Theorem~15.1: the formalization proves the equivalences as a cycle of
separate inequalities with existential constants (the explicit constants
$\theta_1$, $\theta_2$ of the paper are not tracked), it works with the
$L^2$-normalized Haar coefficients, and there is no single wrap-up statement
packaging all four norms at once.}
\begin{lstlisting}[style=lean,basicstyle=\ttfamily\scriptsize]
theorem exists_standardRepresentationNorm_le_const_mul_haarL2RepresentationNorm
    (G : GoodGridSpace (α := α)) [DecidableEq (Set α)]
    (F : UnbalancedHaarWavelet.FullHaarSystem (G := HaarRepresentation.GridOf G))
    [DecidableEq F.Index]
    (s : ℝ) (hs : 0 < s)
    (p : ℝ≥0∞) [Fact (1 ≤ p)] (hp_top : p < ∞)
    (q : ℝ≥0∞) [Fact (1 ≤ q)] :
    ∃ C : ℝ≥0∞, C ≠ ∞ ∧
      ∀ (f : α → ℂ) (hfint : Integrable f G.grid.μ),
        HaarRepresentation.haarL2RepresentationNorm G F s p q f hfint ≠ ∞ →
          standardRepresentationNorm G F s hs p hp_top q f hfint ≠ ∞ ∧
            standardRepresentationNorm G F s hs p hp_top q f hfint ≤
              C * HaarRepresentation.haarL2RepresentationNorm G F s p q f hfint := ...

theorem exists_haarL2RepresentationNorm_le_const_mul_meanOscillationNorm
    (G : GoodGridSpace (α := α)) [DecidableEq (Set α)]
    (F : UnbalancedHaarWavelet.FullHaarSystem (G := GridOf G))
    (s : ℝ) (p : ℝ≥0∞) [Fact (1 ≤ p)] (hp_top : p < ∞)
    (q : ℝ≥0∞) [Fact (1 ≤ q)] :
    ∃ C : ℝ≥0∞, C ≠ ∞ ∧
      ∀ (f : α → ℂ) (hf : Integrable f G.grid.μ), MemLp f p G.grid.μ →
        MeanOscillation.meanOscillationNorm G s p q f ≠ ∞ →
          haarL2RepresentationNorm G F s p q f hf ≤
            C * MeanOscillation.meanOscillationNorm G s p q f := ...

theorem exists_meanOscillationNorm_le_const_mul_souzaBesovNorm
    (G : GoodGridSpace (α := α)) [DecidableEq (Set α)]
    (s : ℝ) (hs : 0 < s)
    (p : ℝ≥0∞) [Fact (1 ≤ p)] (hp_top : p < ∞)
    (q : ℝ≥0∞) [Fact (1 ≤ q)] :
    ∃ C : ℝ≥0∞, C ≠ ∞ ∧
      ∀ (g : SouzaBesovSpace G s p q hs Fact.out (ne_of_lt hp_top))
        (f : α → ℂ),
          f =ᵐ[G.grid.μ] ((g : Lp ℂ p G.toWeakGridSpace.measure) : α → ℂ) →
            meanOscillationNorm G s p q f ≤
              C * ENNReal.ofReal
                (WeakGridSpace.BesovishSpace.Norm_Costpq
                  (souzaAtomFamily G s p hs Fact.out (ne_of_lt hp_top)) q g) := ...

theorem finite_standardRepresentationNorm_implies_memBesov_and_standardRepresentation
    (G : GoodGridSpace (α := α)) [DecidableEq (Set α)]
    (F : UnbalancedHaarWavelet.FullHaarSystem (G := HaarRepresentation.GridOf G))
    [DecidableEq F.Index]
    (s : ℝ) (hs : 0 < s)
    (p : ℝ≥0∞) [Fact (1 ≤ p)] (hp_top : p < ∞)
    (q : ℝ≥0∞) [Fact (1 ≤ q)]
    (f : α → ℂ) (hf : Integrable f G.grid.μ)
    (hN : standardRepresentationNorm G F s hs p hp_top q f hf ≠ ∞) :
    ∃ hfLp : MemLp f p G.grid.μ,
      ∃ g : SouzaBesovSpace G s p q hs Fact.out (ne_of_lt hp_top),
        ∃ R :
          WeakGridSpace.LpGridRepresentation
            (souzaAtomFamily G s p hs Fact.out (ne_of_lt hp_top))
            (g : Lp ℂ p G.toWeakGridSpace.measure),
          (g : Lp ℂ p G.toWeakGridSpace.measure) = hfLp.toLp f ∧
            R.block = canonicalStandardBlockSeq G F s hs p hp_top f hf ∧
              WeakGridSpace.LpGridRepresentation.FinitePQCost (q := q) R ∧
                WeakGridSpace.LpGridRepresentation.pqCostENNReal (q := q) R =
                  standardRepresentationNorm G F s hs p hp_top q f hf ∧
                  WeakGridSpace.LpGridRepresentation.pqCost (q := q) R ≤
                    (standardRepresentationNorm G F s hs p hp_top q f hf).toReal ∧
                    WeakGridSpace.BesovishSpace.Norm_Costpq
                        (souzaAtomFamily G s p hs Fact.out (ne_of_lt hp_top)) q g ≤
                      (standardRepresentationNorm G F s hs p hp_top q f hf).toReal := ...
\end{lstlisting}

\begin{corollary}[Reading the standard coefficients]\label{fou}
For each $P\in \mathcal{P}$ the coefficient $k_P^f$ of the standard atomic
representation of $f$ is a (finite) linear combination of the Haar
coefficients $d^f_S$, $S \in \mathcal H_Q$, of $f$, and the standard
representation of any $f \in \mathcal{B}^s_{p,q}$ converges to $f$ in
$L^p$ with
$$\Big(\sum_k \big( \sum_{P\in \mathcal{P}^k} |k_P^f|^p\big)^{q/p}\Big)^{1/q}\leq C\, |f|_{\mathcal{B}^s_{p,q}}.$$
\end{corollary}

\formalpar{In
\lean{GoodGrid/AlternativeRepresentationsAndNorms/HaarRepresentationNorm.lean},
together with the theorems of the previous \emph{Formalization} paragraph.
The Haar coefficients are the definition \lean{Coeff}, and the theorem
\lean{hasSum\_coeff\_smul\_l2normalizedHaar\_toLp} proves that the
normalized Haar expansion of any $f \in L^\beta$, $1 < \beta < \infty$,
converges to $f$. Paper Corollary~15.2 states moreover that
$f \mapsto k^f_P$ extends to a bounded linear functional on $L^1$; this
extension is not formalized.}
\begin{lstlisting}[style=lean]
theorem hasSum_coeff_smul_l2normalizedHaar_toLp
    (G : GoodGridSpace (α := α)) [DecidableEq (Set α)]
    (F : UnbalancedHaarWavelet.FullHaarSystem (G := GridOf G))
    [DecidableEq F.Index]
    (β : ℝ≥0∞) (hβ_one : 1 < β) (hβ_top : β < ∞)
    (f : α → ℂ) (hf : MemLp f β G.grid.μ) :
    letI : Fact (1 ≤ β) := ⟨le_of_lt hβ_one⟩
    HasSum
      (fun i : F.Index =>
        Coeff G F f
            (by
              letI : IsFiniteMeasure G.grid.μ := (GridOf G).isFinite
              exact hf.integrable (le_of_lt hβ_one))
            i •
          (l2normalizedHaar_memLp G F β i).toLp (L2normalizedHaar G F i))
      (hf.toLp f) := ...
\end{lstlisting}

\section{Positive cone}\label{positivesec}

\filebox{GoodGrid/PositiveCone.lean}{GoodGrid/BesovSpace.lean,
GoodGrid/AlternativeRepresentationsAndNorms/StandarRepresentationNormleqBesovNorm.lean,
WeakGrid/Multipliers.lean}

Positivity is where Souza's atoms shine: since canonical atoms are
\emph{positive multiples of indicators}, a representation with nonnegative
coefficients is a sum of nonnegative functions, and no cancellation can
occur. This rigidity has strong structural consequences, proved in this
file. The support of a positive function is (mod $m$) a countable union of
cells, and --- a sharpening obtained in the formalization --- in \emph{any}
positive representation of a function supported in a cell, all active
cells must sit inside that cell: a coefficient cannot be ``wasted''
outside the support, because nothing can cancel it. This support rigidity
is the engine of the positive-cone non-Archimedean theorems of
Section~\ref{multsec}. The original motivation comes from dynamics: the
fixed points of transfer operators of expanding maps --- invariant
densities --- are nonnegative and often discontinuous, and the positive
cone is their natural home inside the Besov space. The file depends on the
standard-representation comparison of Section~\ref{alt1} (hence its
position in the import order): the decomposition $f = f_+ - f_-$ is
obtained by splitting the coefficients of the \emph{standard}
representation, whose linearity and canonicity make the positive parts
well defined.

We say that $f$ is $\mathcal{B}^s_{p,q}$-positive if there is a
$\mathcal{B}^s_{p,q}$-representation
$$f= \sum_k \sum_{P\in \mathcal{P}^k} c_Pa_P$$
where $c_P\geq 0$ and $a_P$ is the canonical $(s,p)$-Souza's atom supported on
$P$. The set of all $\mathcal{B}^s_{p,q}$-positive functions is a convex cone
in $\mathcal{B}^s_{p,q}$, denoted $\mathcal{B}^{s+}_{p,q}$, with the
``norm''
$$|f|_{\mathcal{B}^{s+}_{p,q}}=\inf \Big( \sum_k \big( \sum_{P\in \mathcal{P}^k} c_P^p \big)^{q/p} \Big)^{1/q},$$
the infimum running over all positive representations of
$f$.

\formalpar{In \lean{GoodGrid/PositiveCone.lean}. A positive representation
(canonical atoms, nonnegative coefficients) is the predicate
\lean{SouzaPositiveRepresentation}, and the cone of
$\mathcal{B}^s_{p,q}$-positive functions is \lean{souzaPositiveCone}, with
the positive ``norm'' \lean{souzaPositiveNorm}. The formalization also uses
the wider notion of a \emph{cone-positive} representation
(\lean{SouzaConePositiveRepresentation}): real coefficients $\geq 0$ and
Souza's atoms which are a.e.\ nonnegative real-valued, not necessarily
canonical. The file also provides the additive infrastructure used in
the non-Archimedean theorems of Section~\ref{multsec}: sums and scalar multiples of positive representations
(\lean{souzaPositiveRepresentationAdd*},
\lean{exists\_souzaPositiveRepresentation\_finset\_sum}) with levelwise
Minkowski control of the cost
(\lean{souzaPositiveRepresentationAdd\_levelCoeffRoot\_le}).}
\begin{lstlisting}[style=lean]
def SouzaPositiveRepresentation
    (G : GoodGridSpace (α := α)) (s : ℝ) (p : ℝ≥0∞)
    (hs : 0 < s) (hp : 1 ≤ p) (hp_top : p ≠ ∞)
    [Fact (1 ≤ p)]
    {f : Lp ℂ p G.toWeakGridSpace.measure}
    (R : WeakGridSpace.LpGridRepresentation
      (souzaAtomFamily G s p hs hp hp_top) f) : Prop :=
  ∀ k, SouzaPositiveLevelBlock G s p hs hp hp_top (R.block k)

def SouzaConePositiveRepresentation
    (G : GoodGridSpace (α := α)) (s : ℝ) (p : ℝ≥0∞)
    (hs : 0 < s) (hp : 1 ≤ p) (hp_top : p ≠ ∞)
    [Fact (1 ≤ p)]
    {f : Lp ℂ p G.toWeakGridSpace.measure}
    (R : WeakGridSpace.LpGridRepresentation
      (souzaAtomFamily G s p hs hp hp_top) f) : Prop :=
  ∀ k, SouzaConePositiveLevelBlock G s p hs hp hp_top (R.block k)
\end{lstlisting}

\begin{proposition}\label{posdec}
Let $0 < s < 1/p$.
\begin{itemize}
\item[A.] If $f\in \mathcal{B}^s_{p,q}$ is (a.e.) real valued then one can
find $f_+,f_-\in \mathcal{B}^{s+}_{p,q}$ such that
$$f=f_+-f_-,$$
with the positive costs of $f_+$ and $f_-$ controlled in terms of
$|f|_{\mathcal{B}^{s}_{p,q}}$; an analogous four-term decomposition holds
for complex valued $f$.
\item[B.] The positive cone $\mathcal{B}^{s+}_{p,q}$ is dense in the cone of
nonnegative functions of $L^p$.
\end{itemize}
\end{proposition}

\formalpar{In \lean{GoodGrid/PositiveCone.lean}. Item~A is
\lean{exists\_souzaPositive\_decomposition\_of\_aeRealValued} (the
decomposition $f = f_+ - f_-$ with controlled positive costs, for a.e.\ real
valued $f$) and
\lean{exists\_souzaPositive\_decomposition\_of\_aeComplexValued} (the
four-term decomposition for complex valued $f$); in the paper this appears
as an unnumbered observation in Section~13, and the complex case is not
stated there. Item~B is
\lean{souzaPositiveCone\_dense\_in\_LpNonnegativeCone}, density of the
positive cone in the nonnegative cone of $L^p$; not stated in the paper.}
\begin{lstlisting}[style=lean,basicstyle=\ttfamily\scriptsize]
theorem exists_souzaPositive_decomposition_of_aeRealValued
    (G : GoodGridSpace (α := α)) (s : ℝ) (p q : ℝ≥0∞)
    (hs : 0 < s) (hp : 1 ≤ p) (hp_top : p ≠ ∞)
    [Fact (1 ≤ p)] [Fact (1 ≤ q)] :
    ∃ C : ℝ,
      0 ≤ C ∧
        ∀ f : WeakGridSpace.BesovishSpace
            (souzaAtomFamily G s p hs hp hp_top) q,
          SouzaBesovAEEqRealValued G s p q hs hp hp_top f →
            ∃ u v : WeakGridSpace.BesovishSpace
                (souzaAtomFamily G s p hs hp hp_top) q,
              SouzaPositiveElement G s p q hs hp hp_top u ∧
                SouzaPositiveElement G s p q hs hp hp_top v ∧
                f = u - v ∧
                souzaPositiveNorm G s p q hs hp hp_top u ≤
                  ENNReal.ofReal C *
                    ENNReal.ofReal (WeakGridSpace.BesovishSpace.Norm_Costpq
                      (souzaAtomFamily G s p hs hp hp_top) q f) ∧
                souzaPositiveNorm G s p q hs hp hp_top v ≤
                  ENNReal.ofReal C *
                    ENNReal.ofReal (WeakGridSpace.BesovishSpace.Norm_Costpq
                      (souzaAtomFamily G s p hs hp hp_top) q f) := ...

theorem souzaPositiveCone_dense_in_LpNonnegativeCone
    (G : GoodGridSpace (α := α)) (s : ℝ) (β q : ℝ≥0∞)
    (hs : 0 < s) (hβ : 1 ≤ β) (hβ_top : β ≠ ∞)
    [Fact (1 ≤ β)] [Fact (1 ≤ q)] :
    LpNonnegativeCone G β ⊆
      closure (SouzaPositiveConeInLbeta G s β q hs hβ hβ_top) := ...
\end{lstlisting}

\begin{proposition}[Support of positive functions]\label{possupp}
If $f \in \mathcal{B}^{s+}_{p,q}$ then its support
$$\operatorname{supp} f = \{ x\in I\colon \ f(x)\neq 0\}$$
is, up to a set of zero measure, a countable union of elements of
$\mathcal{P}$. Moreover, in any positive representation of $f$, the
coefficient of a cell $Q$ vanishes unless $Q$ is contained (mod $m$) in this
union.
\end{proposition}

\formalpar{In \lean{GoodGrid/PositiveCone.lean} (paper Proposition~13.1).
The first sentence is
\lean{support\_ae\_countable\_iUnion\_goodGridCells\_of\_souzaPositiveFunction},
which exhibits the support of a positive function as an a.e.\ countable
union of good-grid cells. The second sentence is the support theorem
\lean{souzaPositiveRepresentation\_coeff\_eq\_zero\_of\_not\_subset\_cell},
a sharper statement obtained in the formalization: in any positive
representation of a function supported in a cell, every active cell is
contained in that cell. A companion ``linchpin'' lemma,
\lean{souzaPositiveRepresentation\_ae\_ne\_zero\_on\_active\_cell}, proves
that a positively represented function is a.e.\ nonzero on every active
cell of its representation; both are used repeatedly in the positive-cone
non-Archimedean arguments of Section~\ref{multsec}.}
\begin{lstlisting}[style=lean]
theorem support_ae_countable_iUnion_goodGridCells_of_souzaPositiveFunction
    (G : GoodGridSpace (α := α)) (s : ℝ) (p q : ℝ≥0∞)
    (hs : 0 < s) (hp : 1 ≤ p) (hp_top : p ≠ ∞)
    [Fact (1 ≤ p)] [Fact (1 ≤ q)]
    {f : α → ℂ}
    (hf : SouzaPositiveFunction G s p q hs hp hp_top f) :
    IsAECountableUnionOfGoodGridCells G {x | f x ≠ 0} := ...

theorem souzaPositiveRepresentation_coeff_eq_zero_of_not_subset_cell
    (G : GoodGridSpace (α := α)) (s : ℝ) (p q : ℝ≥0∞)
    (hs : 0 < s) (hp : 1 ≤ p) (hp_top : p ≠ ∞)
    [Fact (1 ≤ p)] [Fact (1 ≤ q)]
    (Qc : GoodGridCell G)
    {g : Lp ℂ p G.toWeakGridSpace.measure}
    (R : WeakGridSpace.LpGridRepresentation (souzaAtomFamily G s p hs hp hp_top) g)
    (hR : SouzaPositiveRepresentation G s p hs hp hp_top R)
    (hsupp : ∀ᵐ x ∂ G.toWeakGridSpace.measure, x ∉ Qc.cell → (g : α → ℂ) x = 0)
    {n : ℕ} (P : WeakGridSpace.LevelCell G.toWeakGridSpace n)
    (hP : ¬ P.1 ⊆ Qc.cell) :
    (R.block n).coeff P = 0 := ...
\end{lstlisting}

\section{Dirac's approximations}\label{diracsec}

\filebox{GoodGrid/DiracApproximations.lean}{GoodGrid/AlternativeRepresentationsAndNorms/standardRepresentation.lean,
GoodGrid/PositiveCone.lean}

This section documents the formalization of Section~17 of the published
paper. For a cell $Q \in \mathcal{P}^{k_0}$ the \emph{Dirac kernel}
$\mathbb{1}_Q/m(Q)$ is the approximation of the identity at scale $k_0$
attached to the grid. The theme of the section is that the partial sums of
the standard representation of an integrable function $f$
(Section~\ref{alt1}) evaluate, at any point $x_0 \in Q$, to the average of
$f$ against this kernel. Two consequences follow at once: the partial sums
are bounded by the local $L^\infty$ norm of $f$, and --- since at a point
$x_0$ only the unique ancestor cell at each level contributes --- the
\emph{ancestor coefficient sums} of the standard representation are
controlled by $|f|_\infty$. This is exactly the input that the proof of
Pointwise Multipliers~II (Section~\ref{multsec}) needs from this section.
The proof goes through the unbalanced Haar system: a telescoping identity
expands the Dirac kernel in the Haar functions supported on the ancestors
of $Q$, the evaluation identity follows by strong induction on the level,
and the regrouping of the Haar expansion into the standard representation
(Section~\ref{alt1}) transports the result to the standard partial sums.

\begin{proposition}\label{boup}
Let $f \in L^1$. Fix $k_0 \in \mathbb{N}$, a point $x_0$, and let
$Q \in \mathcal{P}^{k_0}$ be the cell containing $x_0$. Write $f_{k_0}$ for
the partial sum, up to scale $k_0$, of the standard representation of $f$
(the Haar partial sum regrouped on the cells, Section~\ref{alt1}). Then:
\begin{itemize}
\item[A.] $|f_{k_0}(x_0)| \leq |f\,\mathbb{1}_Q|_{\infty}$. Moreover, if
$f = \sum_k \sum_{P \in \mathcal{P}^k} c_P a_P$ is \emph{any} positive
$(s,p)$-Souza representation of $f$, then
$$\sum_{J \supseteq Q} c_J\, m(J)^{s-1/p} \leq |f\,\mathbb{1}_Q|_{\infty},$$
the sum running over the ancestors $J$ of $Q$ of level at most $k_0$.
\item[B.] $\displaystyle f_{k_0}(x_0) = \frac{1}{m(Q)} \int_Q f\, dm$.
\end{itemize}
\end{proposition}

\formalpar{In \lean{GoodGrid/DiracApproximations.lean} (paper
Proposition~17.1), namespace \lean{DiracApproximation}. The Dirac kernel is
\lean{diracKernel}, and the telescoping expansion of the kernel in the Haar
system is \lean{diracKernel\_eq\_alpha\_add\_sumDown}, transported from
\lean{Normalized\_indicator\_in\_Haar\_Wavelets\_span} of the dependency
\lean{UnbalancedHaarWavelet}. The partial sums are \lean{partialHaarSum}
and \lean{partialStandardSum}, identified pointwise by
\lean{partialHaarSum\_eq\_partialStandardSum}. The evaluation identity
$f_{k_0}(x_0) = \int f\, \mathbb{1}_Q/m(Q)\, dm$ is
\lean{partialHaarSum\_eq\_integral\_mul\_diracKernel} (proved by strong
induction on the level), and item~B is its standard-representation form
\lean{claimB}. The first bound of item~A is \lean{claimA\_standard}; the
positive-representation bound is \lean{claimA\_positive}, where the
ancestor sum is expressed through \lean{ancestorTerm} (the contribution of
the unique level-$k$ ancestor, with the canonical normalization
$m(J)^{s-1/p}$), and the proof goes through a.e.\ convergence of a
subsequence of the partial sums together with their monotonicity. The
levelwise collapse of the standard blocks at a point to the single ancestor
contribution is \lean{standardLevelSum\_eq\_ancestor\_term}. Compared with
the paper: the partial sums of the standard representation use the
\emph{tilde} atoms of Section~\ref{alt1}, whose value at $x_0$ has modulus
at most $m(J)^{s-1/p}$ but need not equal it; the exact form with
$m(J)^{s-1/p}$, as in item~A, holds for canonical atoms (and is the form
used in Section~\ref{multsec}). Axioms checked: \lean{propext},
\lean{Classical.choice}, \lean{Quot.sound}.}
\begin{lstlisting}[style=lean]
theorem claimA_standard (G : GoodGridSpace (α := α))
    [DecidableEq (Set α)]
    (F : UnbalancedHaarWavelet.FullHaarSystem (G := HaarRepresentation.GridOf G))
    (f : α → ℂ) (hf : Integrable f G.grid.μ)
    (Q : GoodGridCell G) {x₀ : α} (hx₀ : x₀ ∈ Q.cell) :
    ENNReal.ofReal ‖partialHaarSum G F f hf Q.level x₀‖ ≤
      eLpNorm (Set.indicator Q.cell f) ∞ G.grid.μ := ...

theorem claimA_positive (G : GoodGridSpace (α := α))
    (s : ℝ) (p : ℝ≥0∞)
    (hs : 0 < s) (hp : 1 ≤ p) (hp_top : p ≠ ∞) [Fact (1 ≤ p)]
    {f : Lp ℂ p G.toWeakGridSpace.measure}
    (R : WeakGridSpace.LpGridRepresentation (souzaAtomFamily G s p hs hp hp_top) f)
    (hR : SouzaPositiveRepresentation G s p hs hp hp_top R)
    (Q : GoodGridCell G) :
    ENNReal.ofReal (∑ k ∈ Finset.range (Q.level + 1),
        ancestorTerm G s p hs hp hp_top R Q k) ≤
      eLpNorm (Set.indicator Q.cell (⇑f)) ∞ G.toWeakGridSpace.measure := ...

theorem claimB (G : GoodGridSpace (α := α)) [DecidableEq (Set α)]
    (F : UnbalancedHaarWavelet.FullHaarSystem (G := HaarRepresentation.GridOf G))
    (p : ℝ≥0∞) (s : ℝ) (f : α → ℂ) (hf : Integrable f G.grid.μ)
    (Q : GoodGridCell G) {x₀ : α} (hx₀ : x₀ ∈ Q.cell) :
    partialStandardSum G F p s f hf Q.level x₀ =
      ∫ x, f x * ((diracKernel G Q x : ℝ) : ℂ) ∂G.grid.μ := ...
\end{lstlisting}

\section{Pointwise multipliers acting on $\mathcal{B}^s_{p,q}$}\label{multsec}

This section documents the directory \lean{GoodGrid/Multipliers/} (eight
files plus the standalone interface), which formalizes Section~18 of the
published paper; one statement, Pointwise Multipliers~II, is still in
progress. From here on $0 < s < 1/p$. The guiding question is: for
which $g$ is $f \mapsto gf$ bounded on $\mathcal{B}^s_{p,q}$? Because the
space is built from Souza's atoms, and the product of $g$ with a Souza's
atom is just $g$ restricted to a cell and rescaled, the natural test is to
multiply $g$ against every atom --- this gives the \emph{selfs} classes,
and the theory consists in showing that this obviously necessary condition
is, in various quantified senses, also sufficient. The subsections below
follow the import order of the files.

\subsection{The selfs classes and their seminorms}

\filebox{GoodGrid/Multipliers/Definition.lean}{WeakGrid/Multipliers.lean,
GoodGrid/BesovSpace.lean}

Let $g\colon I \rightarrow \mathbb{C}$ be a measurable function. We say that
$g$ is a {\bf pointwise multiplier} acting on $\mathcal{B}^s_{p,q}$ if the
transformation $G(f)=gf$ defines a bounded operator in $\mathcal{B}^s_{p,q}$.
We denote the set of pointwise multipliers by $M(\mathcal{B}^s_{p,q})$, with
the norm
$$|g|_{M(\mathcal{B}^s_{p,q})}=\sup\{|gf|_{\mathcal{B}^s_{p,q}}\colon \ |f|_{\mathcal{B}^s_{p,q}}\leq 1 \}.$$
Define
$$\mathcal{B}^s_{p,q,selfs}=\Big\{ g\in \mathcal{B}^s_{p,q}\colon \ \sup_{a_Q \in \mathcal{A}^{sz}_{s,p}} |ga_Q|_{\mathcal{B}^s_{p,q}} < \infty\Big\},$$
with the norm
$$|g|_{\mathcal{B}^s_{p,q,selfs}}=\sup_{a_Q \in \mathcal{A}^{sz}_{s,p}} |ga_Q|_{\mathcal{B}^s_{p,q}},$$
and, for $t \in \mathbb{N}$, the tail seminorms
$$|g|_{\mathcal{B}^{s,t}_{p,q,selfs}}=\sup_{\substack{ a_Q \in \mathcal{A}^{sz}_{s,p}\\ Q\in \mathcal{P}^j,\ j\geq t}} |ga_Q|_{\mathcal{B}^s_{p,q}}.$$

\formalpar{In \lean{GoodGrid/Multipliers/Definition.lean}. Pointwise
multipliers are the predicate \lean{SouzaPointwiseMultiplier}, the class
$\mathcal{B}^s_{p,q,selfs}$ is \lean{SouzaPointwiseSelfsClass}, and the
tail variant is \lean{SouzaPointwiseSelfsTailClass}: the quantitative form
\lean{SouzaPointwiseSelfsTailBound}~$g$~$C$ records that for every Souza's
atom on a cell of level at least $t$ the product $g\,a_Q$ admits a
representation of cost at most $C$, and the induced seminorm
$|\cdot|_{\mathcal{B}^{s,t}_{p,q,selfs}}$ is
\lean{souzaPointwiseSelfsTailNorm}.}
\begin{lstlisting}[style=lean]
abbrev SouzaPointwiseMultiplier
    (G : GoodGridSpace (α := α)) (s : ℝ) (p q : ℝ≥0∞)
    (hs : 0 < s) (hp : 1 ≤ p) (hp_top : p ≠ ∞)
    [Fact (1 ≤ p)] [Fact (1 ≤ q)]
    (m : α → ℂ) : Prop :=
  WeakGridSpace.IsPointwiseMultiplier
    (A := souzaAtomFamily G s p hs hp hp_top) q m

abbrev SouzaPointwiseSelfsClass
    (G : GoodGridSpace (α := α)) (s : ℝ) (p q : ℝ≥0∞)
    (hs : 0 < s) (hp : 1 ≤ p) (hp_top : p ≠ ∞)
    [Fact (1 ≤ p)] [Fact (1 ≤ q)]
    (m : α → ℂ) : Prop :=
  WeakGridSpace.PointwiseSelfsClass
    (A := souzaAtomFamily G s p hs hp hp_top) q m
\end{lstlisting}

\subsection{Restriction to a cell}

\filebox{GoodGrid/Multipliers/Besovspq.lean}{GoodGrid/Multipliers/Definition.lean,
WeakGrid/Transmutation.lean}

The first multiplier theorem is for the simplest candidates: indicators of
cells. The proof is another instance of transmutation --- restricting a
Souza's atom $a_Q$ to a cell $W$ produces either $0$, the atom itself, or
(when $W \subset Q$) a rescaled atom on $W$, with geometric decay in the
level gap; Proposition~8.1(A) and (B) of the paper then assemble these
local expansions with uniform cost. The lemma is stated on the induced
grid of $W$, and the uniformity of $C_{22}$ in $W$ is what the
non-Archimedean arguments below need.

\begin{lemma}[Restriction to a cell]\label{incint3}
Let $W\in \mathcal{P}$. The restriction map
$$r\colon \mathcal{B}^s_{p,q}(I, \mathcal{P}, \mathcal{A}^{sz}_{s,p})\rightarrow \mathcal{B}^s_{p,q}(W, \mathcal{P}_W, \mathcal{A}^{sz}_{s,p}),
\qquad r(f)=1_Wf,$$
is continuous: there is $C_{10} \geq 1$, independent of $W$, such that
\begin{itemize}
\item[A.] For every $f \in \mathcal{B}^s_{p,q}$ we have
$$|1_Wf|_{\mathcal{B}^s_{p,q}(W, \mathcal{P}_W, \mathcal{A}^{sz}_{s,p})}\leq C_{10}\, |f|_{\mathcal{B}^s_{p,q}}.$$
In particular $1_W\in M(\mathcal{B}^s_{p,q})$ with
$|1_W|_{M(\mathcal{B}^s_{p,q})} \le C_{10}$, uniformly in $W$.
\item[B.] For every $\mathcal{B}^{s+}_{p,q}$-representation
$f=\sum_{k}\sum_{Q} c_Qa_Q$ one can find a positive representation
$$1_Wf= \sum_{k}\sum_{Q} d_Q\,a_Q$$
on the induced grid with cost at most
$C_{10}$ times the cost of the original one, and such that $d_Q\neq 0$
implies $Q\subset \operatorname{supp} f$ (mod $m$).
\end{itemize}
\end{lemma}

\formalpar{In \lean{GoodGrid/Multipliers/Besovspq.lean} (paper Lemma~18.2).
The restriction map with its induced-grid representation is
\lean{souzaCellIndicatorRestrictsToInduced}; the cost bounds (uniform in
$W$) are the lemmas
\lean{souzaIndicatorRestrictionBound\_of\_ambientRestrictionTransmutation\_*}
and \lean{souzaCellIndicatorPointwiseMultiplierBound\_uniform}, the latter
giving the explicit multiplier bound $2C+1$ with
$C = {}$\lean{souzaAmbientRestrictionMultiplierConstant}; and the theorem
\lean{souzaCellIndicatorPointwiseMultiplier} concludes that $1_W$ is a
pointwise multiplier of $\mathcal{B}^s_{p,q}$. The paper's strongly regular
domains (Definition~18.7 and Proposition~18.8), of which the present lemma
is the grid-cell case, are not formalized in general; the multiplier
statement for $1_W$, $W \in \mathcal{P}$, with a bound uniform in $W$, is
the formalized special case.}
\begin{lstlisting}[style=lean]
theorem souzaCellIndicatorPointwiseMultiplier
    (G : GoodGridSpace (α := α)) (s : ℝ) (p q : ℝ≥0∞)
    (hs : 0 < s) (hp : 1 ≤ p) (hp_top : p ≠ ∞)
    [Fact (1 ≤ p)] [Fact (1 ≤ q)]
    (hs_lt_inv : s < (p.toReal)⁻¹)
    (W : GoodGridCell G) :
    SouzaPointwiseMultiplier G s p q hs hp hp_top
      (W.cell.indicator fun _ => (1 : ℂ)) := ...

theorem souzaCellIndicatorPointwiseMultiplierBound_uniform
    (G : GoodGridSpace (α := α)) (s : ℝ) (p q : ℝ≥0∞)
    (hs : 0 < s) (hp : 1 ≤ p) (hp_top : p ≠ ∞)
    [Fact (1 ≤ p)] [Fact (1 ≤ q)]
    (hs_lt_inv : s < (p.toReal)⁻¹)
    (W : GoodGridCell G) :
    SouzaPointwiseMultiplierBound G s p q hs hp hp_top
      (W.cell.indicator fun _ => (1 : ℂ))
      (2 * souzaAmbientRestrictionMultiplierConstant G s p + 1) := ...
\end{lstlisting}

\subsection{The multipliers of $\mathcal{B}^s_{1,1}$}

\filebox{GoodGrid/Multipliers/Besovs11.lean}{GoodGrid/Multipliers/Definition.lean}

For $p = q = 1$ the multiplier problem has a complete and clean answer:
the necessary condition is sufficient. The reason is that for an
$\ell^1$-type cost one can substitute the representation of $f$ into $gf$
atom by atom, $gf = \sum c_Q (g a_Q)$, and the costs of the pieces
$g a_Q$ simply add up --- no interaction between cells survives. This is
paper Proposition~18.1.

\begin{proposition}\label{m11}
We have that
$$M(\mathcal{B}^s_{1,1})=\mathcal{B}^s_{1,1,selfs}.$$
\end{proposition}

\formalpar{This is paper Proposition~18.1. The equivalence is the theorem
\lean{souzaPointwiseMultiplier\_iff\_souzaPointwiseSelfsClass\_one\_one}:
the inclusion $\mathcal{B}^s_{1,1,selfs}\subset M(\mathcal{B}^s_{1,1})$ is
proved in \lean{GoodGrid/Multipliers/Besovs11.lean}, while the converse,
which holds for all $p,q$, is
\lean{souzaPointwiseSelfsClass\_of\_souzaPointwiseMultiplier} in
\lean{GoodGrid/Multipliers/Besovspq.lean}.}
\begin{lstlisting}[style=lean]
theorem souzaPointwiseMultiplier_iff_souzaPointwiseSelfsClass_one_one
    (G : GoodGridSpace (α := α)) (s : ℝ)
    (hs : 0 < s) {m : α → ℂ} :
    SouzaPointwiseMultiplier G s (1 : ℝ≥0∞) (1 : ℝ≥0∞)
      hs le_rfl ENNReal.one_ne_top m ↔
    SouzaPointwiseSelfsClass G s (1 : ℝ≥0∞) (1 : ℝ≥0∞)
      hs le_rfl ENNReal.one_ne_top m := ...
\end{lstlisting}

\subsection{$selfs$ multipliers are bounded}

\filebox{GoodGrid/Multipliers/MultipliersareBounded.lean}{GoodGrid/BesovAtoms.lean,
GoodGrid/Multipliers/Definition.lean, GoodGrid/Multipliers/Besovspq.lean}

A function in a selfs class is automatically bounded. The proof is a
pretty martingale argument (paper Proposition~18.3): testing $g$ against
the canonical atom of a cell $Q$ and applying the $L^p$ embedding of
Proposition~\ref{lp} to the product bounds the $L^p$ average of $g$ over
$Q$, uniformly over all cells; since the cells generate the
$\sigma$-algebra ($G_8$), L\'evy's upward theorem lets these averages
converge to $|g(x)|$ at almost every point, giving the sup bound. The
formalization sharpens the paper in a direction the non-Archimedean
theorems require: the constant is \emph{independent of the tail level
$t$}, so that multipliers known only on deep cells are still uniformly
bounded.

\begin{proposition}[$selfs$ classes are bounded]\label{selfsbdd}
There is $C_{11}=C_{11}(\lambda_1,\lambda_2,C_1,s,p,q)$ such that for every
$t \in \mathbb{N}$ and every $g$ with
$|g|_{\mathcal{B}^{s,t}_{p,q,selfs}}<\infty$,
$$|g(x)| \leq C_{11}\, |g|_{\mathcal{B}^{s,t}_{p,q,selfs}}
\quad \text{for } m\text{-a.e. } x.$$
In particular $\mathcal{B}^s_{p,q,selfs} \subset L^\infty$ and this inclusion
is continuous.
\end{proposition}

\formalpar{In \lean{GoodGrid/Multipliers/MultipliersareBounded.lean}. The
pointwise estimate from a single tail bound is
\lean{souzaPointwiseSelfsTailBound\_norm\_ae\_le}; the estimate in terms of
the tail seminorm, with the explicit constant displayed in the code below,
is \lean{souzaPointwiseSelfsTailNorm\_norm\_ae\_le}; and the qualitative
$L^\infty$ membership of the whole class is
\lean{souzaPointwiseSelfsTailClass\_norm\_ae\_bounded}. This sharpens paper
Proposition~18.3 in two respects: the estimate is proved for the tail
seminorms $|\cdot|_{\mathcal{B}^{s,t}_{p,q,selfs}}$ with a constant
\emph{independent of $t$} (this uniformity is what the non-Archimedean
arguments below require), and the formalization also provides a
\emph{canonical} variant (\lean{SouzaPointwiseCanonicalSelfsTailBound}), in
which $g$ is only tested against products with canonical Souza's atoms,
together with the bridges \lean{Bound.toCanonical} and
\lean{SouzaPositivePointwiseSelfsTailBound.toCanonical}; the canonical
variant is what the infinite positive-cone theorem below consumes.}
\begin{lstlisting}[style=lean]
theorem souzaPointwiseSelfsTailNorm_norm_ae_le
    (G : GoodGridSpace (α := α))
    (s : ℝ) (p q : ℝ≥0∞)
    (hs : 0 < s) (hp : 1 ≤ p) (hp_top : p ≠ ∞)
    [Fact (1 ≤ p)] [Fact (1 ≤ q)]
    (hs_lt_inv : s < (p.toReal)⁻¹)
    {t : ℕ} {m : α → ℂ}
    (hm : SouzaPointwiseSelfsTailClass G s p q hs hp hp_top t m) :
    ∀ᵐ x ∂G.grid.μ,
      ‖m x‖ ≤
        souzaBesovLpLocalEmbeddingConstant G s p q *
          (2 * souzaAmbientRestrictionMultiplierConstant G s p + 1) *
            souzaPointwiseSelfsTailNorm G s p q hs hp hp_top t m := ...
\end{lstlisting}

\subsection{Non-Archimedean behaviour in $\mathcal{B}^\beta_{p,\tilde{q},selfs}$}

\filebox{GoodGrid/Multipliers/NonArchimedeanProperty.lean (statements
exposed by NonArchimedeanPropertyPositiveStandalone.lean)}{GoodGrid/Multipliers/Definition.lean,
GoodGrid/BesovAtoms.lean, GoodGrid/Multipliers/Besovspq.lean,
GoodGrid/Multipliers/MultipliersareBounded.lean, GoodGrid/PositiveCone.lean,
WeakGrid/Scales.lean}

These are the deepest theorems of the repository (paper Proposition~18.4
and Remark~18.5). The phenomenon: if the supports of the multipliers
$g_i$ are \emph{separated by the grid} --- every cell carrying a nonzero
coefficient of $f$ meets essentially one support, which is what
hypothesis A quantifies --- then each atom of the representation of $f$
``sees'' only one $g_i$, and the cost of $(\sum_i g_i)f$ is governed by
the \emph{largest} relevant multiplier norm instead of their sum: the
estimate behaves like an ultrametric bound, whence \emph{non-Archimedean}.
The level hypothesis B makes the separation scale-respecting: an atom may
only meet a multiplier whose tail level has already been reached. The
proof multiplies the representation of $f$ atomwise, expands each product
$g_i a_Q$ by the restriction lemma on the induced grid of $Q$, transmutes
the resulting $\beta$-smooth representations down to smoothness $s$
(transmutation again --- here \lean{WeakGrid/Scales.lean} provides the
$\beta \to s$ rescaling), and reassembles; the bounded overlap of the
supports is exactly what keeps the reassembled cost additive in the right
way. The positive versions add the support and positivity conclusions of
Remark~18.5, whose proofs rest on the support rigidity of
Section~\ref{positivesec}; the infinite-index versions require, in
addition, the pointwise theory: absolute convergence of $\sum_i g_i(x)$
on $\{f \neq 0\}$, with the limit identified through the compactness of
representations of Section~\ref{complsec}.

If $g_i \in M(\mathcal{B}^s_{p,q})$, the naive estimate is
$$\Big|\big(\sum_i g_i\big)f\Big|_{\mathcal{B}^s_{p,q}}\leq \Big(\sum_i
|g_i|_{M(\mathcal{B}^s_{p,q})}\Big)\,|f|_{\mathcal{B}^s_{p,q}}.$$
Remarkably, when the supports of the $g_i$ are separated by the grid, one
gets a far better estimate.

\begin{theorem}[Non-Archimedean property, finite version]\label{sepa}
Let $\beta > s$ and $\tilde q \in [1,\infty]$. There is
$C_{12}=C_{12}(\lambda_1,\lambda_2,C_1,s,\beta,p,q,\tilde q)$ with the
following property. Let $\Lambda$ be finite, $g_i$ with
$|g_i|_{\mathcal{B}^{\beta,t_i}_{p,\tilde{q},selfs}} < \infty$ for $i\in
\Lambda$ and $t_i \in \mathbb{N}$. Consider a function $f$ with a
$\mathcal{B}^s_{p,q}$-representation
$$ f=\sum_k \sum_{Q\in \mathcal{P}^k} c_Q a_Q$$
by canonical Souza's atoms satisfying
\begin{itemize}
\item[A.] We have
$$\sup_{\substack{ Q\in \mathcal{P} \\ c_Q\neq 0}} \
\sum_{i\colon\, Q\cap \operatorname{supp} g_i \neq \emptyset} |g_i|_{ \mathcal{B}^{\beta,t_i}_{p,\tilde{q},selfs}} \leq N.$$
\item[B.] If $Q\in \mathcal{P}^k$ satisfies $c_Q\neq 0$ and
$Q\cap \operatorname{supp} g_i \neq \emptyset$ then $k\geq t_i$.
\end{itemize}
Then $(\sum_{i\in\Lambda} g_i)f \in \mathcal{B}^s_{p,q}$ and we can find a
$\mathcal{B}^s_{p,q}$-representation
\begin{equation} \label{rre} \big(\sum_{i\in\Lambda} g_i\big)f= \sum_k \sum_{Q\in \mathcal{P}^k} d_Q a_Q\end{equation}
by canonical Souza's atoms such that
\begin{equation}\label{reep33} \Big(\sum_k \big( \sum_{Q\in \mathcal{P}^k} |d_Q|^p \big)^{q/p} \Big)^{1/q}
\leq C_{12}\, N\, \Big(\sum_k \big( \sum_{Q\in \mathcal{P}^k} |c_Q|^p \big)^{q/p} \Big)^{1/q}.
\end{equation}
\end{theorem}

\formalpar{The theorem \lean{souzaNonArchimedeanPropertyLambdaFinite} in
\lean{GoodGrid/Multipliers/NonArchimedeanProperty.lean} (paper
Proposition~18.4 for finite $\Lambda$) produces, from the bounded-overlap
hypothesis A and the level/scale hypothesis B, a representation of the
product $(\sum_i g_i)f$ of cost at most $C_{12}\,N$ times the cost of the
representation of $f$, with $C_{12}$ independent of the multipliers, of
$N$, and of $\Lambda$. The representations of $f$ and of the product are by
\emph{canonical} Souza's atoms (\lean{SouzaCanonicalRepresentation}), which
is no loss of generality for Souza's atoms and is the precise form
formalized.}
\begin{lstlisting}[style=lean,basicstyle=\ttfamily\scriptsize]
theorem souzaNonArchimedeanPropertyLambdaFinite
    (G : GoodGridSpace (α := α))
    (s β : ℝ) (p q qtilde : ℝ≥0∞)
    (hs : 0 < s) (hβ : 0 < β) (hβs : s < β)
    (hβ_lt_inv : β < (p.toReal)⁻¹)
    (hp : 1 ≤ p) (hp_top : p ≠ ∞)
    [Fact (1 ≤ p)] [Fact (1 ≤ q)] [Fact (1 ≤ qtilde)] :
    ∃ Cgen : ℝ,
      0 ≤ Cgen ∧
      ∀ (Λ : Finset ℕ) (t : ℕ → ℕ) (g : ℕ → α → ℂ) (N : ℝ)
        (f : α → ℂ)
        (x : WeakGridSpace.BesovishSpace
          (souzaAtomFamily G s p hs hp hp_top) q)
        (R : WeakGridSpace.LpGridRepresentation
          (souzaAtomFamily G s p hs hp hp_top)
          (x : Lp ℂ p G.toWeakGridSpace.measure)),
          0 ≤ N →
          WeakGridSpace.RepresentsFunction
            (G := G.toWeakGridSpace) (p := p) f
            (x : Lp ℂ p G.toWeakGridSpace.measure) →
          WeakGridSpace.LpGridRepresentation.FinitePQCost (q := q) R →
          (∀ i ∈ Λ,
            ∃ C : ℝ,
              SouzaPointwiseSelfsTailBound
                G β p qtilde hβ hp hp_top (t i) (g i) C) →
          (∀ k (Q : WeakGridSpace.LevelCell G.toWeakGridSpace k),
            (R.block k).coeff Q ≠ 0 →
              nonArchimedeanRelevantTailSelfsSum
                G β p qtilde hβ hp hp_top Λ t g Q ≤ N) →
          (∀ k (Q : WeakGridSpace.LevelCell G.toWeakGridSpace k) i,
            i ∈ Λ →
              (R.block k).coeff Q ≠ 0 →
                goodGridLevelCellMeetsSupport G Q (g i) →
                  t i ≤ k) →
          ∃ y : WeakGridSpace.BesovishSpace
              (souzaAtomFamily G s p hs hp hp_top) q,
            ∃ S : WeakGridSpace.LpGridRepresentation
                (souzaAtomFamily G s p hs hp hp_top)
                (y : Lp ℂ p G.toWeakGridSpace.measure),
              WeakGridSpace.RepresentsFunction
                (G := G.toWeakGridSpace) (p := p)
                (fun z => (∑ i ∈ Λ, g i z) * f z)
                (y : Lp ℂ p G.toWeakGridSpace.measure) ∧
              WeakGridSpace.LpGridRepresentation.pqCost (q := q) S ≤
                Cgen * N *
                  WeakGridSpace.LpGridRepresentation.pqCost (q := q) R := ...
\end{lstlisting}

\begin{theorem}[Non-Archimedean property, infinite version]\label{sepainf}
In the setting of Theorem~\ref{sepa} let $\Lambda \subset \mathbb{N}$ be
infinite, and replace hypothesis A.\ by the summability condition
$$\sum_{i \in \Lambda} T_i \leq N, \qquad
T_i := \sup_{\substack{ Q\in \mathcal{P},\ c_Q\neq 0 \\ Q\cap \operatorname{supp} g_i \neq \emptyset}}
|g_i|_{ \mathcal{B}^{\beta,t_i}_{p,\tilde{q},selfs}}\ ,$$
understood as a series with values in $[0,\infty]$. Then, for $m$-a.e.\ $x$
in $\{f \neq 0\}$, the series $\sum_{i\in\Lambda} g_i(x)$ converges
absolutely, with
$$\sum_{i\in\Lambda} |g_i(x)| \leq C_{12}\, N;$$
the function
$h = (\sum_{i\in\Lambda} g_i) f$ (defined a.e.\ as this pointwise limit
times $f$) belongs to $L^p$ and to $\mathcal{B}^s_{p,q}$, and it has a
canonical representation (\ref{rre}) satisfying
(\ref{reep33}).
\end{theorem}

\formalpar{The theorem \lean{souzaNonArchimedeanProperty} in
\lean{GoodGrid/Multipliers/NonArchimedeanProperty.lean} proves all the
conclusions above: pointwise absolute summability of $\sum_i g_i(x)$ with
bound $C_{12} N$ a.e.\ on $\{f \neq 0\}$, existence of the limit function
$h$ with $\|h\| \leq C_{12} N \|f\|$ a.e., membership of $h$ in $L^p$, and a
canonical representation of $h$ with the cost bound (\ref{reep33}). The
pointwise part is \lean{exists\_nonArchimedeanInfinite\_pointwise\_hasSum};
the identification of the limit representation goes through the compactness
of representations with uniform cost
(\lean{exists\_strongly\_convergent\_subseq\_of\_uniform\_pqCost}) applied
to the truncations to finite initial segments of $\Lambda$. Paper
Proposition~18.4 treats finite and infinite $\Lambda$ together, reducing to
the finite case; the formalization keeps them separate because the infinite
case has genuinely more content.}
\begin{lstlisting}[style=lean,basicstyle=\ttfamily\scriptsize]
theorem souzaNonArchimedeanProperty
    (G : GoodGridSpace (α := α))
    (s β : ℝ) (p q qtilde : ℝ≥0∞)
    (hs : 0 < s) (hβ : 0 < β) (hβs : s < β)
    (hβ_lt_inv : β < (p.toReal)⁻¹)
    (hp : 1 ≤ p) (hp_top : p ≠ ∞)
    [Fact (1 ≤ p)] [Fact (1 ≤ q)] [Fact (1 ≤ qtilde)] :
    ∃ Cgen : ℝ,
      0 ≤ Cgen ∧
      1 ≤ Cgen ∧
      ∀ (Λ : Set ℕ) (t : ℕ → ℕ) (g : ℕ → α → ℂ) (N : ℝ)
        (f : α → ℂ)
        (x : WeakGridSpace.BesovishSpace
          (souzaAtomFamily G s p hs hp hp_top) q)
        (R : WeakGridSpace.LpGridRepresentation
          (souzaAtomFamily G s p hs hp hp_top)
          (x : Lp ℂ p G.toWeakGridSpace.measure)),
          0 ≤ N →
          WeakGridSpace.RepresentsFunction
            (G := G.toWeakGridSpace) (p := p) f
            (x : Lp ℂ p G.toWeakGridSpace.measure) →
          WeakGridSpace.LpGridRepresentation.FinitePQCost (q := q) R →
          (∀ i ∈ Λ,
            ∃ C : ℝ,
              SouzaPointwiseSelfsTailBound
                G β p qtilde hβ hp hp_top (t i) (g i) C) →
          (∀ k (Q : WeakGridSpace.LevelCell G.toWeakGridSpace k),
            (R.block k).coeff Q ≠ 0 →
              ∃ T : ℝ,
                HasSum
                  (fun i : {i // i ∈ Λ} =>
                    nonArchimedeanRelevantTailSelfsInfiniteTerm
                      G β p qtilde hβ hp hp_top Λ t g Q i)
                  T ∧
                T ≤ N) →
          (∀ k (Q : WeakGridSpace.LevelCell G.toWeakGridSpace k) i,
            i ∈ Λ →
              (R.block k).coeff Q ≠ 0 →
                goodGridLevelCellMeetsSupport G Q (g i) →
                  t i ≤ k) →
          ∃ h : α → ℂ,
            ∃ absSum : α → ℝ,
              (∀ᵐ z ∂G.toWeakGridSpace.measure,
                f z ≠ 0 →
                  HasSum
                    (fun i : {i // i ∈ Λ} => ‖g i.1 z‖)
                    (absSum z) ∧
                  absSum z ≤ Cgen * N) ∧
              (∀ᵐ z ∂G.toWeakGridSpace.measure,
                HasSum
                  (fun i : {i // i ∈ Λ} => g i.1 z * f z)
                  (h z)) ∧
              (∀ᵐ z ∂G.toWeakGridSpace.measure,
                ‖h z‖ ≤ Cgen * N * ‖f z‖) ∧
              (∃ hmem : MemLp h p G.toWeakGridSpace.measure,
                ‖MemLp.toLp h hmem‖ ≤
                  Cgen * N * ‖(x : Lp ℂ p G.toWeakGridSpace.measure)‖) ∧
              ∃ y : WeakGridSpace.BesovishSpace
                  (souzaAtomFamily G s p hs hp hp_top) q,
                ∃ S : WeakGridSpace.LpGridRepresentation
                    (souzaAtomFamily G s p hs hp hp_top)
                    (y : Lp ℂ p G.toWeakGridSpace.measure),
                  WeakGridSpace.RepresentsFunction
                    (G := G.toWeakGridSpace) (p := p) h
                    (y : Lp ℂ p G.toWeakGridSpace.measure) ∧
                  WeakGridSpace.LpGridRepresentation.FinitePQCost (q := q) S ∧
                  WeakGridSpace.LpGridRepresentation.pqCost (q := q) S ≤
                    Cgen * N *
                      WeakGridSpace.LpGridRepresentation.pqCost (q := q) R := ...
\end{lstlisting}

\begin{theorem}[Non-Archimedean property, positive version]\label{posrem}
In the setting of Theorems~\ref{sepa} and~\ref{sepainf} (finite or infinite
$\Lambda$, with the respective hypotheses, the tail seminorms being replaced
by their positive-cone counterparts
$$|g|_{\mathcal{B}^{\beta+,t}_{p,\tilde q,selfs}}
=\sup_{\substack{ a_Q \in \mathcal{A}^{sz}_{\beta,p}\\ Q\in \mathcal{P}^j,\ j\geq t}} |ga_Q|_{\mathcal{B}^{\beta+}_{p,\tilde q}}$$
for the functions $g_i$), assume that each $g_i$ is
$\mathcal{B}^{\beta,t_i}_{p,\tilde q}$-positive. Let $R$ be the given
canonical representation $(c_Q)_Q$ of $f$. Then the representation
$S=(d_Q)_Q$ of $(\sum_i g_i)f$ produced by
Theorems~\ref{sepa}/\ref{sepainf} can be chosen so that, in addition:
\begin{itemize}
\item[i.] ({\it positivity}) if $R$ is a positive representation
($c_Q \geq 0$ for all $Q$), then $S$ is a cone-positive representation: all
$d_Q$ are nonnegative reals and the corresponding Souza's atoms are a.e.\
nonnegative;
\item[ii.] ({\it support}) unconditionally --- without assuming $c_Q\geq 0$
--- every cell $Q$ active in $S$ ($d_Q \neq 0$) satisfies: there exists
$i \in \Lambda$ with $g_i \neq 0$ a.e.\ on $Q$.
\end{itemize}
\end{theorem}

\formalpar{In the file
\lean{NonArchimedeanPropertyPositiveStandalone.lean} of
\lean{GoodGrid/Multipliers/}, with cores in
\lean{NonArchimedeanProperty.lean} and
\lean{GoodGrid/PositiveCone.lean}. The finite case is the theorem
\lean{souzaNonArchimedeanPropertyPositiveCone}: it produces the
representation $S$ of the product with the cost bound of
Theorem~\ref{sepa} and, in addition, the support conclusion ii.\ and the
positivity implication i. The infinite case is
\lean{souzaNonArchimedeanPropertyPositiveConeInfinite}
($\Lambda \subset \mathbb{N}$ infinite, with condition A as an
$[0,\infty]$-valued series bound, no summability witness needed): it proves
all the conclusions of Theorem~\ref{sepainf} plus i.\ and ii.; the limit
arguments pass ii.\ through the convergence of the coefficients and i.\
through the closedness of the nonnegative real ray in $\mathbb{C}$ and
a.e.\ convergence of a subsequence of the atoms. The auxiliary product
representation behind both is
\lean{exists\_nonArchimedeanProductRepresentation\_positive}, whose key
local ingredient
(\lean{exists\_nonArchimedeanLocalTransmutationData\_pos}) assembles, one
multiplier at a time, positive local transmutation data on the induced
grids. Compared with paper Remark~18.5 the statement was made precise as
follows: (a) the hypothesis on the representation $R$ of $f$ is only that
it is canonical (arbitrary complex coefficients); (b) the two conclusions
are separated according to their true strengths --- the support conclusion
ii.\ holds unconditionally, and is weakened to its a.e.\ form ($g_i \neq 0$
a.e.\ on $Q$ rather than $Q \subset \operatorname{supp} g_i$); (c) the
positivity conclusion i.\ produces a \emph{cone-positive} representation
(nonnegative real coefficients, a.e.\ nonnegative Souza's atoms), not
necessarily one by canonical atoms. Axioms of all four non-Archimedean
theorems checked: \lean{propext}, \lean{Classical.choice},
\lean{Quot.sound}.}
\begin{lstlisting}[style=lean,basicstyle=\ttfamily\scriptsize]
theorem souzaNonArchimedeanPropertyPositiveCone
    (G : GoodGridSpace (α := α))
    (s β : ℝ) (p q qtilde : ℝ≥0∞)
    (hs : 0 < s) (hβ : 0 < β) (hβs : s < β)
    (hβ_lt_inv : β < (p.toReal)⁻¹)
    (hp : 1 ≤ p) (hp_top : p ≠ ∞)
    [Fact (1 ≤ p)] [Fact (1 ≤ q)] [Fact (1 ≤ qtilde)] :
    ∃ Cgen : ℝ,
      0 ≤ Cgen ∧
      ∀ (Λ : Finset ℕ) (t : ℕ → ℕ) (g : ℕ → α → ℂ) (N : ℝ)
        (f : α → ℂ)
        (x : WeakGridSpace.BesovishSpace
          (souzaAtomFamily G s p hs hp hp_top) q)
        (R : WeakGridSpace.LpGridRepresentation
          (souzaAtomFamily G s p hs hp hp_top)
          (x : Lp ℂ p G.toWeakGridSpace.measure)),
          0 ≤ N →
          WeakGridSpace.RepresentsFunction
            (G := G.toWeakGridSpace) (p := p) f
            (x : Lp ℂ p G.toWeakGridSpace.measure) →
          WeakGridSpace.LpGridRepresentation.FinitePQCost (q := q) R →
          SouzaCanonicalRepresentation G s p hs hp hp_top R →
          (∀ i ∈ Λ,
            SouzaPositiveFunction G β p qtilde hβ hp hp_top (g i)) →
          (∀ i ∈ Λ,
            ∃ C : ℝ≥0∞,
              SouzaPositivePointwiseSelfsTailBound
                G β p qtilde hβ hp hp_top (t i) (g i) C) →
          (∀ k (Q : WeakGridSpace.LevelCell G.toWeakGridSpace k),
            (R.block k).coeff Q ≠ 0 →
              nonArchimedeanRelevantPositiveTailSelfsSum
                G β p qtilde hβ hp hp_top Λ t g Q ≤ ENNReal.ofReal N) →
          (∀ k (Q : WeakGridSpace.LevelCell G.toWeakGridSpace k) i,
            i ∈ Λ →
              (R.block k).coeff Q ≠ 0 →
                goodGridLevelCellMeetsSupport G Q (g i) →
                  t i ≤ k) →
          ∃ y : WeakGridSpace.BesovishSpace
              (souzaAtomFamily G s p hs hp hp_top) q,
            ∃ S : WeakGridSpace.LpGridRepresentation
                (souzaAtomFamily G s p hs hp hp_top)
                (y : Lp ℂ p G.toWeakGridSpace.measure),
              WeakGridSpace.RepresentsFunction
                (G := G.toWeakGridSpace) (p := p)
                (fun z => (∑ i ∈ Λ, g i z) * f z)
                (y : Lp ℂ p G.toWeakGridSpace.measure) ∧
              WeakGridSpace.LpGridRepresentation.FinitePQCost (q := q) S ∧
              WeakGridSpace.LpGridRepresentation.pqCost (q := q) S ≤
                Cgen * N *
                  WeakGridSpace.LpGridRepresentation.pqCost (q := q) R ∧
              -- [ii] support: holds with only canonical atoms on `R` (no sign on `c_Q`).
              (∀ k (Q : WeakGridSpace.LevelCell G.toWeakGridSpace k),
                (S.block k).coeff Q ≠ 0 →
                  ∃ i ∈ Λ, ∀ᵐ z ∂(G.toWeakGridSpace.measure.restrict Q.1), g i z ≠ 0) ∧
              -- [i] positivity: only when `R` is fully positive (canonical + `c_Q ≥ 0`).
              (SouzaPositiveRepresentation G s p hs hp hp_top R →
                SouzaConePositiveRepresentation G s p hs hp hp_top S) := ...

theorem souzaNonArchimedeanPropertyPositiveConeInfinite
    (G : GoodGridSpace (α := α))
    (s β : ℝ) (p q qtilde : ℝ≥0∞)
    (hs : 0 < s) (hβ : 0 < β) (hβs : s < β)
    (hβ_lt_inv : β < (p.toReal)⁻¹)
    (hp : 1 ≤ p) (hp_top : p ≠ ∞)
    [Fact (1 ≤ p)] [Fact (1 ≤ q)] [Fact (1 ≤ qtilde)] :
    ∃ Cgen : ℝ,
      0 ≤ Cgen ∧
      1 ≤ Cgen ∧
      ∀ (Λ : Set ℕ) (t : ℕ → ℕ) (g : ℕ → α → ℂ) (N : ℝ)
        (f : α → ℂ)
        (x : WeakGridSpace.BesovishSpace
          (souzaAtomFamily G s p hs hp hp_top) q)
        (R : WeakGridSpace.LpGridRepresentation
          (souzaAtomFamily G s p hs hp hp_top)
          (x : Lp ℂ p G.toWeakGridSpace.measure)),
          0 ≤ N →
          WeakGridSpace.RepresentsFunction
            (G := G.toWeakGridSpace) (p := p) f
            (x : Lp ℂ p G.toWeakGridSpace.measure) →
          WeakGridSpace.LpGridRepresentation.FinitePQCost (q := q) R →
          SouzaCanonicalRepresentation G s p hs hp hp_top R →
          (∀ i ∈ Λ,
            SouzaPositiveFunction G β p qtilde hβ hp hp_top (g i)) →
          (∀ i ∈ Λ,
            ∃ C : ℝ≥0∞,
              SouzaPositivePointwiseSelfsTailBound
                G β p qtilde hβ hp hp_top (t i) (g i) C) →
          (∀ k (Q : WeakGridSpace.LevelCell G.toWeakGridSpace k),
            (R.block k).coeff Q ≠ 0 →
              (∑' i : {i // i ∈ Λ},
                nonArchimedeanRelevantPositiveTailSelfsInfiniteTerm
                  G β p qtilde hβ hp hp_top Λ t g Q i) ≤ ENNReal.ofReal N) →
          (∀ k (Q : WeakGridSpace.LevelCell G.toWeakGridSpace k) i,
            i ∈ Λ →
              (R.block k).coeff Q ≠ 0 →
                goodGridLevelCellMeetsSupport G Q (g i) →
                  t i ≤ k) →
          ∃ h : α → ℂ,
            ∃ absSum : α → ℝ,
              (∀ᵐ z ∂G.toWeakGridSpace.measure,
                f z ≠ 0 →
                  HasSum
                    (fun i : {i // i ∈ Λ} => ‖g i.1 z‖)
                    (absSum z) ∧
                  absSum z ≤ Cgen * N) ∧
              (∀ᵐ z ∂G.toWeakGridSpace.measure,
                HasSum
                  (fun i : {i // i ∈ Λ} => g i.1 z * f z)
                  (h z)) ∧
              (∀ᵐ z ∂G.toWeakGridSpace.measure,
                ‖h z‖ ≤ Cgen * N * ‖f z‖) ∧
              (∃ hmem : MemLp h p G.toWeakGridSpace.measure,
                ‖MemLp.toLp h hmem‖ ≤
                  Cgen * N * ‖(x : Lp ℂ p G.toWeakGridSpace.measure)‖) ∧
              ∃ y : WeakGridSpace.BesovishSpace
                  (souzaAtomFamily G s p hs hp hp_top) q,
                ∃ S : WeakGridSpace.LpGridRepresentation
                    (souzaAtomFamily G s p hs hp hp_top)
                    (y : Lp ℂ p G.toWeakGridSpace.measure),
                  WeakGridSpace.RepresentsFunction
                    (G := G.toWeakGridSpace) (p := p) h
                    (y : Lp ℂ p G.toWeakGridSpace.measure) ∧
                  WeakGridSpace.LpGridRepresentation.FinitePQCost (q := q) S ∧
                  WeakGridSpace.LpGridRepresentation.pqCost (q := q) S ≤
                    Cgen * N *
                      WeakGridSpace.LpGridRepresentation.pqCost (q := q) R ∧
                  (∀ k (Q : WeakGridSpace.LevelCell G.toWeakGridSpace k),
                    (S.block k).coeff Q ≠ 0 →
                      ∃ i ∈ Λ,
                        ∀ᵐ z ∂(G.toWeakGridSpace.measure.restrict Q.1),
                          g i z ≠ 0) ∧
                  (SouzaPositiveRepresentation G s p hs hp hp_top R →
                    SouzaConePositiveRepresentation G s p hs hp hp_top S) := ...
\end{lstlisting}

\subsection{The tail selfs classes consist of multipliers}

\filebox{GoodGrid/Multipliers/SelfsSubsetMultipliers.lean}{GoodGrid/Multipliers/NonArchimedeanProperty.lean}

The first consequence of the non-Archimedean property is the continuous
inclusion of the tail \lean{selfs} classes in the multiplier space.

\begin{corollary}\label{selfsmult}
Let $0 < s < \beta < 1/p$ and $q, \tilde q \in [1,\infty]$, and fix a cutoff
level $t \in \mathbb{N}$. Then
$$\mathcal{B}^{\beta,t}_{p,\tilde q,selfs} \subset M(\mathcal{B}^s_{p,q}),$$
and the inclusion is continuous:
$|g|_{M(\mathcal{B}^s_{p,q})} \leq C\, |g|_{\mathcal{B}^{\beta,t}_{p,\tilde q,selfs}}$
with a constant $C$ independent of $g$.
\end{corollary}

\formalpar{In \lean{GoodGrid/Multipliers/SelfsSubsetMultipliers.lean}
(paper Corollary~18.6). The quantitative form is
\lean{exists\_souzaSelfsMultiplierConstant}: a tail \lean{selfs} bound $C$
for $g$ produces the multiplier operator bound $C_{mult}\cdot C$. The
inclusion is
\lean{souzaPointwiseMultiplier\_of\_souzaPointwiseSelfsTailClass} and the
continuity at the level of (semi)norms is
\lean{souzaPointwiseMultiplierNorm\_le\_const\_mul\_selfsTailNorm}. The
proof applies the non-Archimedean property (Theorem~\ref{sepa}) to the
one-element family $\{g\}$, after the level-lowering step
\lean{souzaPointwiseSelfsTailBound\_levelZero}: a tail \lean{selfs} bound
from level $t$ onward yields a full (level-$0$) bound with constant
multiplied by the number of level-$t$ cells, because an atom on a shallow
cell is the finite disjoint sum of its level-$t$ pieces, each a scaled atom
on a tail cell. Axioms checked: \lean{propext}, \lean{Classical.choice},
\lean{Quot.sound}.}
\begin{lstlisting}[style=lean]
theorem exists_souzaSelfsMultiplierConstant
    (G : GoodGridSpace (α := α)) (s β : ℝ) (p q qtilde : ℝ≥0∞) (t : ℕ)
    (hs : 0 < s) (hβ : 0 < β) (hβs : s < β)
    (hβ_lt_inv : β < (p.toReal)⁻¹)
    (hp : 1 ≤ p) (hp_top : p ≠ ∞)
    [Fact (1 ≤ p)] [Fact (1 ≤ q)] [Fact (1 ≤ qtilde)] :
    ∃ Cmult : ℝ,
      0 ≤ Cmult ∧
      ∀ (g : α → ℂ) (C : ℝ),
        SouzaPointwiseSelfsTailBound G β p qtilde hβ hp hp_top t g C →
        SouzaPointwiseMultiplierBound G s p q hs hp hp_top g (Cmult * C) := ...

theorem souzaPointwiseMultiplierNorm_le_const_mul_selfsTailNorm
    (G : GoodGridSpace (α := α)) (s β : ℝ) (p q qtilde : ℝ≥0∞) (t : ℕ)
    (hs : 0 < s) (hβ : 0 < β) (hβs : s < β)
    (hβ_lt_inv : β < (p.toReal)⁻¹)
    (hp : 1 ≤ p) (hp_top : p ≠ ∞)
    [Fact (1 ≤ p)] [Fact (1 ≤ q)] [Fact (1 ≤ qtilde)] :
    ∃ Cmult : ℝ,
      0 ≤ Cmult ∧
      ∀ g : α → ℂ,
        SouzaPointwiseSelfsTailClass G β p qtilde hβ hp hp_top t g →
        WeakGridSpace.pointwiseMultiplierNorm
            (A := souzaAtomFamily G s p hs hp hp_top) q g ≤
          Cmult * souzaPointwiseSelfsTailNorm G β p qtilde hβ hp hp_top t g := ...
\end{lstlisting}

\subsection{Strongly regular domains and Pointwise Multipliers I}

\filebox{GoodGrid/Multipliers/StronglyRegularDomains.lean}{GoodGrid/Multipliers/NonArchimedeanPropertyPositiveStandalone.lean}

The positive-cone version of the non-Archimedean property was designed for
indicators, and strongly regular domains are the class of sets whose
indicators it handles.

\begin{definition}\label{srd}
Let $a \in \mathbb{R}$, $K \geq 0$ and $k_1 \in \mathbb{N}$. A measurable
set $\Omega$ is an $(a, K, k_1)$-\textbf{strongly regular domain} if for
every grid cell $Q$ of level at least $k_1$ the intersection $Q \cap
\Omega$ is an exactly disjoint countable union of grid cells,
$$Q \cap \Omega = \bigcup_k \bigcup_{P \in \mathcal{F}^k} P,$$
with $\mathcal{F}^k \subset \mathcal{P}^k$ and the level-by-level cost
control
$$\sum_{P \in \mathcal{F}^k} m(P)^a \leq K\, m(Q)^a
\quad \text{for every } k.$$
\end{definition}

\begin{proposition}\label{pos2}
If $\Omega$ is a $(1-\beta p, K, k_1)$-strongly regular domain then
$$|\mathbb{1}_\Omega|_{\mathcal{B}^{\beta+,k_1}_{p,\infty,selfs}} \leq K^{1/p}:$$
for every grid cell $Q$ of level at least $k_1$, the product of
$\mathbb{1}_\Omega$ with the canonical $(\beta,p)$-Souza's atom on $Q$ admits
a positive Souza representation of positive $(p,\infty)$-gauge at most
$K^{1/p}$. In particular $\mathbb{1}_\Omega \in \mathcal{B}^{\beta+}_{p,\infty}$.
\end{proposition}

\formalpar{In \lean{GoodGrid/Multipliers/StronglyRegularDomains.lean}
(paper 18.7--18.8). The decomposition data inside one cell is the structure
\lean{StronglyRegularDecomposition} (families $\mathcal{F}^k$ of genuine
level-$k$ cells, exact disjoint cover of $Q \cap \Omega$, disjointness also
across levels, and the cost bound), and the definition is
\lean{StronglyRegularDomain}. Proposition~\ref{pos2} is
\lean{souzaPositiveSelfsTailBound\_of\_stronglyRegularDomain}, with the
weighted variant
\lean{souzaPositiveSelfsTailBound\_smul\_of\_stronglyRegularDomain}
($\Theta \cdot \mathbb{1}_\Omega$ with constant $\Theta\, K^{1/p}$) and the
membership \lean{souzaPositiveFunction\_of\_stronglyRegularDomain}: the
positive representation of $\mathbb{1}_\Omega \cdot a_Q$ is read off
directly from the decomposition families.}
\begin{lstlisting}[style=lean]
def StronglyRegularDomain
    (G : GoodGridSpace (α := α)) (Ω : Set α) (a K : ℝ) (k₁ : ℕ) : Prop :=
  MeasurableSet Ω ∧
    ∀ Q : GoodGridCell G, k₁ ≤ Q.level →
      Nonempty (StronglyRegularDecomposition G Ω a K Q)

theorem souzaPositiveSelfsTailBound_of_stronglyRegularDomain
    (G : GoodGridSpace (α := α)) (β : ℝ) (p : ℝ≥0∞)
    (hβ : 0 < β) (hp : 1 ≤ p) (hp_top : p ≠ ∞)
    [Fact (1 ≤ p)]
    {Ω : Set α} {K : ℝ} {k₁ : ℕ}
    (hΩ : StronglyRegularDomain G Ω (1 - β * p.toReal) K k₁) :
    SouzaPositivePointwiseSelfsTailBound G β p ∞ hβ hp hp_top k₁
      (Ω.indicator fun _ => (1 : ℂ))
      (ENNReal.ofReal (K ^ (p.toReal)⁻¹)) := ...
\end{lstlisting}

\begin{theorem}[Pointwise Multipliers I]\label{pm1}
There is a constant $C$ with the following property. Let $\Omega_i$, for
$i$ in a finite set $\Lambda$, be $(1-\beta p, K_i, t_i)$-strongly regular
domains, let $\Theta_i > 0$ be weights, and let $f$ be represented by a
finite-cost canonical Souza representation $R = (c_Q)_Q$ in
$\mathcal{B}^s_{p,q}$ such that
\begin{itemize}
\item[A.] on every active cell $Q$ ($c_Q \neq 0$), the total weighted cost
$\sum_{i\colon\, Q \cap \Omega_i \neq \emptyset} \Theta_i\, K_i^{1/p}$ is at
most $N$;
\item[B.] every active cell meeting $\Omega_i$ has level at least $t_i$.
\end{itemize}
Then $g = \sum_{i \in \Lambda} \Theta_i\, \mathbb{1}_{\Omega_i}$ is a
multiplier of the represented function: the product $g f$ has a Souza
representation $S$ with finite cost and
$$\mathrm{cost}(S) \leq C\, N\, \mathrm{cost}(R),$$
and moreover (i) every active cell of $S$ is, up to measure zero, contained
in some $\Omega_i$, and (ii) if $R$ is positive then $S$ is cone-positive.
An infinite-family version ($\Lambda \subset \mathbb{N}$ countable) holds
with condition~A stated as an $[0,\infty]$-valued series bound.
\end{theorem}

\formalpar{In \lean{GoodGrid/Multipliers/StronglyRegularDomains.lean}
(paper Proposition~18.9). The finite case is
\lean{souzaPointwiseMultipliersI} and the countable case is
\lean{souzaPointwiseMultipliersIInfinite}; the weighted overlap costs of
condition~A are \lean{stronglyRegularOverlapCost} and
\lean{stronglyRegularOverlapCostInfinite}. Both are direct applications of
the positive-cone non-Archimedean theorems (Theorem~\ref{posrem}) to the
weighted indicators, whose positive tail \lean{selfs} bounds come from
Proposition~\ref{pos2}; the support conclusion ii.\ of Theorem~\ref{posrem}
(nonvanishing a.e.) here upgrades to a.e.\ \emph{containment} in some
$\Omega_i$, since the multipliers are indicators. Axioms checked:
\lean{propext}, \lean{Classical.choice}, \lean{Quot.sound}.}
\begin{lstlisting}[style=lean,basicstyle=\ttfamily\scriptsize]
theorem souzaPointwiseMultipliersI
    (G : GoodGridSpace (α := α))
    (s β : ℝ) (p q : ℝ≥0∞)
    (hs : 0 < s) (hβ : 0 < β) (hβs : s < β)
    (hβ_lt_inv : β < (p.toReal)⁻¹)
    (hp : 1 ≤ p) (hp_top : p ≠ ∞)
    [Fact (1 ≤ p)] [Fact (1 ≤ q)] :
    ∃ Cgen2 : ℝ,
      0 ≤ Cgen2 ∧
      ∀ (Λ : Finset ℕ) (t : ℕ → ℕ) (Ω : ℕ → Set α) (K Θ : ℕ → ℝ) (N : ℝ)
        (f : α → ℂ)
        (x : WeakGridSpace.BesovishSpace
          (souzaAtomFamily G s p hs hp hp_top) q)
        (R : WeakGridSpace.LpGridRepresentation
          (souzaAtomFamily G s p hs hp hp_top)
          (x : Lp ℂ p G.toWeakGridSpace.measure)),
        0 ≤ N →
        (∀ i ∈ Λ,
          StronglyRegularDomain G (Ω i) (1 - β * p.toReal) (K i) (t i)) →
        (∀ i ∈ Λ, 0 < Θ i) →
        WeakGridSpace.RepresentsFunction
          (G := G.toWeakGridSpace) (p := p) f
          (x : Lp ℂ p G.toWeakGridSpace.measure) →
        WeakGridSpace.LpGridRepresentation.FinitePQCost (q := q) R →
        SouzaCanonicalRepresentation G s p hs hp hp_top R →
        (∀ k (Q : WeakGridSpace.LevelCell G.toWeakGridSpace k),
          (R.block k).coeff Q ≠ 0 →
            stronglyRegularOverlapCost G p Λ Ω K Θ Q ≤ N) →
        (∀ k (Q : WeakGridSpace.LevelCell G.toWeakGridSpace k) i,
          i ∈ Λ →
            (R.block k).coeff Q ≠ 0 →
              (Q.1 ∩ Ω i).Nonempty →
                t i ≤ k) →
        ∃ y : WeakGridSpace.BesovishSpace
            (souzaAtomFamily G s p hs hp hp_top) q,
          ∃ S : WeakGridSpace.LpGridRepresentation
              (souzaAtomFamily G s p hs hp hp_top)
              (y : Lp ℂ p G.toWeakGridSpace.measure),
            WeakGridSpace.RepresentsFunction
              (G := G.toWeakGridSpace) (p := p)
              (fun z => (∑ i ∈ Λ,
                (Θ i : ℂ) * (Ω i).indicator (fun _ => (1 : ℂ)) z) * f z)
              (y : Lp ℂ p G.toWeakGridSpace.measure) ∧
            WeakGridSpace.LpGridRepresentation.FinitePQCost (q := q) S ∧
            WeakGridSpace.LpGridRepresentation.pqCost (q := q) S ≤
              Cgen2 * N *
                WeakGridSpace.LpGridRepresentation.pqCost (q := q) R ∧
            (∀ k (Q : WeakGridSpace.LevelCell G.toWeakGridSpace k),
              (S.block k).coeff Q ≠ 0 →
                ∃ i ∈ Λ,
                  ∀ᵐ z ∂(G.toWeakGridSpace.measure.restrict Q.1),
                    z ∈ Ω i) ∧
            (SouzaPositiveRepresentation G s p hs hp hp_top R →
              SouzaConePositiveRepresentation G s p hs hp hp_top S) := ...
\end{lstlisting}

\subsection{$\mathcal{B}^{1/p}_{p,\infty} \cap L^\infty$ and Pointwise
Multipliers II (in progress)}\label{pm2sub}

\filebox{GoodGrid/Multipliers/Bp1overpinftyisMultiplier.lean}{GoodGrid/Multipliers/NonArchimedeanProperty.lean,
GoodGrid/DiracApproximations.lean,
GoodGrid/AlternativeRepresentationsAndNorms/FiniteStandardNormimpliesBesov.lean}

\begin{proposition}[Pointwise Multipliers II]\label{pm2}
Let $0 < s < 1/p$. There is a constant $C$ such that every
$g \in \mathcal{B}^{1/p}_{p,\infty} \cap L^\infty$ is a pointwise
multiplier of $\mathcal{B}^s_{p,q}$, with
$$|g|_{M(\mathcal{B}^s_{p,q})} \leq C\, |g|_{\mathcal{B}^{1/p}_{p,\infty}}
+ |g|_\infty.$$
\end{proposition}

\formalpar{In \lean{GoodGrid/Multipliers/Bp1overpinftyisMultiplier.lean}
(paper Proposition~18.10). \emph{This is the only statement of this
document whose formalization is still in progress.} The statement is
\lean{souzaPointwiseMultipliersII}, with the hypothesis
$g \in \mathcal{B}^{1/p}_{p,\infty} \cap L^\infty$ rendered as: $g$ is
represented by an element \lean{xg} of the $(1/p, p, \infty)$ Souza space
and $\|g\| \leq M$ a.e. Its proof is organized around two inner sublemmas.
The first, \lean{exists\_fouRepresentation}, is \emph{fully proved}: it
packages the input from Corollary \lean{fou} (Section~\ref{alt1}) and
Proposition~17.1.B (Section~\ref{diracsec}) --- the canonical standard
representation of $g$ has $(p,\infty)$-cost controlled by
$|g|_{\mathcal{B}^{1/p}_{p,\infty}}$, and all its ancestor-tower
coefficient sums (\lean{ancestorCoeffSum}) are bounded by $M$, because for
$s = 1/p$ the canonical atoms have value $m(J)^{1/p-1/p} = 1$ on their
cells, so the tower sums are exactly the partial standard sums, i.e.\
averages of $g$ over cells (axioms checked: \lean{propext},
\lean{Classical.choice}, \lean{Quot.sound}). The second sublemma,
\lean{exists\_mult\_product\_representation} --- the paper's splitting
$gf = u_1 + u_2$ of the product into the convolution piece ($J
\subsetneq Q$, geometric kernel $\lambda_2^{(j-k)(1/p-s)}$, factor
$|g|_{\mathcal{B}^{1/p}_{p,\infty}}/(1-\lambda_2^{1/p-s})$) and the
tower-sum piece ($Q \subseteq J$, factor $|g|_\infty$), identified with
$gf$ through $L^1$-convergent truncations and Corollary \lean{compa1}
(Section~\ref{complsec}) --- is the \emph{only declared \lean{sorry} of
the repository}; the blocks of both pieces (\lean{multU1Block},
\lean{multU2Block}) and their coefficient estimates
(\lean{multU1Block\_coeff\_norm\_le},
\lean{multU2Block\_levelCoeffPower\_le}) are already proved. The outer
proof of \lean{souzaPointwiseMultipliersII} --- the
$\varepsilon$-optimization over near-optimal representations of $f$ and
the uniqueness of the product representative in $L^p$ --- is complete.
Until the remaining sublemma is closed, this file is \emph{not} imported
by the aggregate root.}
\begin{lstlisting}[style=lean,basicstyle=\ttfamily\scriptsize]
theorem souzaPointwiseMultipliersII
    (G : GoodGridSpace (α := α)) (s : ℝ) (p q : ℝ≥0∞)
    (hs : 0 < s) (hs_lt_inv : s < (p.toReal)⁻¹)
    (h1p : 0 < (p.toReal)⁻¹) (hp : 1 ≤ p) (hp_top : p ≠ ∞)
    [Fact (1 ≤ p)] [Fact (1 ≤ q)] :
    ∃ Cmult : ℝ,
      0 ≤ Cmult ∧
      ∀ (g : α → ℂ) (M : ℝ)
        (xg : WeakGridSpace.BesovishSpace
          (souzaAtomFamily G (p.toReal)⁻¹ p h1p hp hp_top) ∞),
        WeakGridSpace.RepresentsFunction
          (G := G.toWeakGridSpace) (p := p) g
          (xg : Lp ℂ p G.toWeakGridSpace.measure) →
        (∀ᵐ z ∂G.toWeakGridSpace.measure, ‖g z‖ ≤ M) →
        SouzaPointwiseMultiplierBound G s p q hs hp hp_top g
          (Cmult * WeakGridSpace.BesovishSpace.Norm_Costpq
              (souzaAtomFamily G (p.toReal)⁻¹ p h1p hp hp_top) ∞ xg
            + M) := ...

theorem exists_fouRepresentation
    (G : GoodGridSpace (α := α)) (p : ℝ≥0∞)
    (h1p : 0 < (p.toReal)⁻¹) (hp : 1 ≤ p) (hp_top : p ≠ ∞)
    [Fact (1 ≤ p)] :
    ∃ Cfou : ℝ,
      0 ≤ Cfou ∧
      ∀ (g : α → ℂ) (M : ℝ)
        (xg : WeakGridSpace.BesovishSpace
          (souzaAtomFamily G (p.toReal)⁻¹ p h1p hp hp_top) ∞),
        WeakGridSpace.RepresentsFunction
          (G := G.toWeakGridSpace) (p := p) g
          (xg : Lp ℂ p G.toWeakGridSpace.measure) →
        (∀ᵐ z ∂G.toWeakGridSpace.measure, ‖g z‖ ≤ M) →
        ∃ Rg : WeakGridSpace.LpGridRepresentation
            (souzaAtomFamily G (p.toReal)⁻¹ p h1p hp hp_top)
            (xg : Lp ℂ p G.toWeakGridSpace.measure),
          SouzaCanonicalRepresentation G (p.toReal)⁻¹ p h1p hp hp_top Rg ∧
          WeakGridSpace.LpGridRepresentation.FinitePQCost (q := ∞) Rg ∧
          WeakGridSpace.LpGridRepresentation.pqCost (q := ∞) Rg ≤
            Cfou * WeakGridSpace.BesovishSpace.Norm_Costpq
              (souzaAtomFamily G (p.toReal)⁻¹ p h1p hp hp_top) ∞ xg ∧
          ∀ (k : ℕ) (Q : WeakGridSpace.LevelCell G.toWeakGridSpace k),
            ‖ancestorCoeffSum G Rg Q‖ ≤ M := ...
\end{lstlisting}

\vspace{8mm}
\centerline{\Large\bf III. GUIDE TO THE REPOSITORY}
\vspace{6mm}

\section{Guide to the repository files}\label{guide}

All paths below are relative to the library root
\lean{BesovSpacesGoodGrid/}; the aggregate file
\lean{BesovSpacesGoodGrid.lean} is the main entry point and imports the
whole pipeline.

\subsection{Weak-grid layer (Part~I)}

\begin{itemize}
  \item \lean{WeakGrid/Definition.lean}: weak-grid structures
    (\lean{WeakGrid}, \lean{WeakGridSpace}) and the overlap estimates of
    Section~\ref{partition}.
  \item \lean{WeakGrid/Atoms.lean}: weak-grid cells, local spaces, and atom
    families (Section~\ref{secatoms}).
  \item \lean{WeakGrid/BesovishSpaces.lean}: level blocks, representations,
    costs, Besov-ish spaces, the cost gauge \lean{Norm\_Costpq}, the $L^t$
    embedding and the normed structure (Section~\ref{besovish}).
  \item \lean{WeakGrid/Scales.lean}: scaled atom families and the
    smoothness-scale inclusion (Section~\ref{escala}).
  \item \lean{WeakGrid/Completeness.lean}: representation limits,
    compactness of cost balls, and completeness (Section~\ref{complsec}).
  \item \lean{WeakGrid/InducedGrid.lean}: induced weak grids on a fixed cell
    and the contraction into the ambient Besov-ish space
    (Section~\ref{inducedsec}).
  \item \lean{WeakGrid/Transmutation.lean}: transmutation objects and
    Claims I, II, III, including the explicit Claim~C embedding constants
    (Section~\ref{transsec}).
  \item \lean{WeakGrid/Multipliers.lean}: the abstract multiplier and
    selfs-class predicates on weak grids (Section~\ref{infra}).
  \item \lean{Sums.lean}: block-index types (\lean{BlockIndex},
    \lean{ChildIndex}, \lean{ParentIndex}) and block-sum notation used
    throughout (Section~\ref{infra}).
\end{itemize}

\subsection{Good-grid layer (Part~II)}

\begin{itemize}
  \item \lean{GoodGrid/Definition.lean}: good-grid structures
    (Section~\ref{goodgrids}).
  \item \lean{GoodGrid/BesovSpace.lean}: Souza's atoms and good-grid Besov
    spaces, with completeness, compactness, and density
    (Sections~\ref{secatom} and~\ref{secbesov}).
  \item \lean{GoodGrid/PositiveCone.lean}: the positive cone, positive and
    cone-positive representations, decompositions, density, and the support
    theorems (Section~\ref{positivesec}).
  \item \lean{GoodGrid/DiracApproximations.lean}: the grid Dirac kernels,
    the evaluation of partial standard sums as cell averages, and the
    $L^\infty$ bounds on ancestor coefficient sums of paper
    Proposition~17.1 (Section~\ref{diracsec}).
  \item \lean{GoodGrid/BesovAtoms.lean}: Besov's atoms on good grids and the
    Souza/Besov sandwich theorem (Section~\ref{alt2}).
  \item \lean{GoodGrid/AlternativeRepresentationsAndNorms.lean} and
    \lean{GoodGrid/AlternativeRepresentationsAndNorms/}: the norm-comparison
    layer of Section~\ref{alt1} ---
    \lean{HaarRepresentationNorm.lean} (normalized Haar coefficients and the
    Haar gauge), \lean{HaarParametrizedRepresentation.lean} and
    \lean{ComparingHaarRepresentationsl.lean} (Haar expansions as
    parametrized representations and their comparison),
    \lean{standardRepresentation.lean} (standard atomic
    coefficients and gauge),
    \lean{standardNormleqHaarRepresenstionNorm.lean} (Haar controls
    standard), \lean{StandarRepresentationNormleqBesovNorm.lean} (the Besov
    norm controls the standard norm), \lean{FiniteStandardNormimpliesBesov.lean}
    (finite standard norm endpoint), \lean{FiniteHaarNormimpliesLp.lean}
    (finite Haar norm endpoint), \lean{MeanOscillationNorm.lean}
    (mean-oscillation definitions and local oscillation lemmas),
    \lean{OscillationNormleqBesovNorm.lean} (Besov controls mean
    oscillation), and \lean{HaarNormleqOscillationNorm.lean} (mean
    oscillation controls Haar).
  \item \lean{GoodGrid/Distribution.lean}: cell indicators
    (\lean{atomIndicator}), the test-function module
    (\lean{TestFunctions}), and its algebraic dual
    (\lean{Distributions}).
\end{itemize}

\subsection{Multiplier layer (Part~II, Section~\ref{multsec})}

\begin{itemize}
  \item \lean{GoodGrid/Multipliers.lean} and \lean{GoodGrid/Multipliers/}:
    the pointwise-multiplier layer of Section~\ref{multsec} ---
    \lean{Definition.lean} (the \lean{selfs} classes and seminorms),
    \lean{Besovspq.lean} (induced-cell representations, restriction to a
    cell, and the general inclusion
    $M(\mathcal{B}^s_{p,q}) \subset \mathcal{B}^s_{p,q,selfs}$),
    \lean{Besovs11.lean} (the $\mathcal{B}^s_{1,1}$ characterization),
    \lean{MultipliersareBounded.lean} (boundedness of \lean{selfs}
    multipliers, uniform in the tail level, and the canonical tail
    classes), \lean{NonArchimedeanProperty.lean} (the finite and infinite
    non-Archimedean estimates and the cores of their positive versions),
    \lean{NonArchimedeanPropertyPositiveStandalone.lean} (the
    user-facing positive-cone statements, finite and infinite),
    \lean{SelfsSubsetMultipliers.lean} (the continuous inclusion of the
    tail \lean{selfs} classes in the multiplier space, paper
    Corollary~18.6), \lean{StronglyRegularDomains.lean} (strongly regular
    domains and Pointwise Multipliers~I, paper 18.7--18.9), and
    \lean{Bp1overpinftyisMultiplier.lean} (Pointwise Multipliers~II, paper
    Proposition~18.10; in progress --- it contains the only declared
    \lean{sorry} of the repository and is not yet imported by the
    aggregate root).
\end{itemize}

\subsection{External dependencies}

The repository depends on the external Lean repositories
\lean{UnbalancedHaarWavelet} (grids and unbalanced Haar systems),
\lean{UnconditionalSchauderBasis},
\lean{LaminarFamiliesMaximalBinaryTrees} and \lean{Burkholder} (all by
SmaniaD), which provide the Girardi--Sweldens construction and the
unconditional basis property used in Section~\ref{haar}.

\subsection{Repository status}

At the date of this document (June 11, 2026) a full \lean{lake build}
succeeds (3457 jobs). The project Lean files contain no \lean{admit} and no
project-local \lean{axiom} or \lean{constant} declarations, and exactly one
declared \lean{sorry}: the inner sublemma
\lean{exists\_mult\_product\_representation} (the $u_1+u_2$ construction of
Pointwise Multipliers~II, Section~\ref{pm2sub}) of
\lean{GoodGrid/Multipliers/Bp1overpinftyisMultiplier.lean}, a file not yet
imported by the aggregate root \lean{BesovSpacesGoodGrid.lean}; every
module imported by the root compiles with no \lean{sorry}. The axioms of
the main theorems --- the finite and infinite, non-positive and positive
non-Archimedean properties and their consequences, namely
\begin{center}
\small
\lean{souzaNonArchimedeanProperty},\quad
\lean{souzaNonArchimedeanPropertyLambdaFinite},\\
\lean{souzaNonArchimedeanPropertyPositiveCone},\\
\lean{souzaNonArchimedeanPropertyPositiveConeInfinite},\\
\lean{exists\_souzaSelfsMultiplierConstant},\quad
\lean{souzaPointwiseMultipliersI},\\
\lean{souzaPointwiseMultipliersIInfinite},\quad
\lean{claimA\_standard},\quad
\lean{claimA\_positive},\\
\lean{claimB},\quad
\lean{exists\_fouRepresentation}
\end{center}
--- were checked
with \lean{\#print axioms}: only the standard axioms \lean{propext},
\lean{Classical.choice}, \lean{Quot.sound} appear (no \lean{sorryAx}). The
running verification log is kept in \lean{status.md} at the repository
root.

\begin{thebibliography}{9}

\bibitem{paper} D. Smania,
\emph{Besov-ish spaces through atomic decomposition},
Analysis \& PDE \textbf{15} (2022), no.~1, 123--174.
A prerelease version of the published paper is included in this repository
as \lean{preprint-source/besovish-2020-07-18.tex} (with PDF); it is the
source from which the present document was produced.

\bibitem{gw} M. Girardi and W. Sweldens,
\emph{A new class of unbalanced Haar wavelets that form an unconditional basis for $L_p$ on general measure spaces},
J. Fourier Anal. Appl. \textbf{3} (1997), 457--474.

\bibitem{souzao1} G. S. de Souza,
\emph{Spaces formed by special atoms},
PhD thesis, SUNY Albany, 1980.

\bibitem{repo} D. Smania,
\emph{BesovSpacesGoodGrid: a Lean 4 formalization},
git repository, 2026. Dependencies: \texttt{UnbalancedHaarWavelet},
\texttt{UnconditionalSchauderBasis}, \texttt{LaminarFamiliesMaximalBinaryTrees},
\texttt{Burkholder} (SmaniaD).

\end{thebibliography}

\end{document}
